\documentclass[conference,compsoc]{IEEEtran}
\usepackage[utf8]{inputenc}
\usepackage{xspace}
\usepackage{graphicx}
\usepackage{float}
\usepackage{cuted}
\usepackage{subcaption}
\usepackage{pgfplots}
\usepackage{amsthm}
\usepackage{amsmath}
\usepackage{amssymb}
\usepackage{booktabs}
\usepackage[nocompress]{cite}
\usepackage[colorlinks=false, pdfborder={0 0 0}, hidelinks]{hyperref}
\pgfplotsset{compat=1.16}
\usepgfplotslibrary[colorbrewer]
\usetikzlibrary{pgfplots.colorbrewer}
\usetikzlibrary{calc}
\definecolor{color1}{HTML}{0868ac}
\definecolor{color2}{HTML}{43a2ca}
\definecolor{color3}{HTML}{7bccc4}
\definecolor{color4}{HTML}{a8ddb5}
\definecolor{color5}{HTML}{ccebc5}
\definecolor{myPurple}{HTML}{7B3294}
\definecolor{myOrange}{HTML}{E66100}
\definecolor{myGreen}{HTML}{1B9E77}
\definecolor{myRed}{HTML}{D7191C}

\newtheorem{proposition}{Proposition}

\newtheorem{definition}{Definition}

\begin{document}
\sloppy

\title{Biometric Authentication Mechanisms}

% \author{\IEEEauthorblockN{Dor Malka}
% \IEEEauthorblockA{Technion,\\
% Email: dm29@campus.technion.ac.il}
% \and
% \IEEEauthorblockN{Ittay Eyal}
% \IEEEauthorblockA{Technion,\\
% Email: ittay@technion.ac.il}}

% \renewcommand{\shortauthors}{Malka et al.}
\maketitle
\tableofcontents

\begin{abstract}
Modern authentication commonly uses \emph{biometric credentials} such as fingerprints and faces: 
a sensor evaluates each scan and accepts if its score exceeds a \emph{match threshold}. 
A large body of work improves sensor accuracy but sets this threshold heuristically. 
Recent work studies \emph{authentication mechanisms} that combine credentials to maximize the probability of accepting the user and rejecting the attacker. 
But it treats every credential as \emph{atomic}: presented correctly or not, like a password. 

We model authentication mechanisms with both biometric and atomic credentials. 
For a single-attempt mechanism that uses a biometric credential alone, threshold optimization yields only marginal improvements. 
However, when biometrics are combined with other credentials or used across multiple attempts, threshold optimization becomes consequential.
We prove that a mechanism combining an atomic credential with an optimal-threshold biometric one is better than a mechanism using only the atomic credential; this holds even if the biometric credential is extremely inaccurate. 
Next, we consider mechanisms allowing multiple attempts, as fingerprint readers commonly do. 
We find that each attempt's threshold should be different. 
Using public biometric datasets, we find that optimal thresholds can increase the success probability of multiple-attempt mechanisms by~72 percentage points compared to the heuristic threshold. 

Yet our survey of commercial sensors finds that almost none allow setting a specific threshold. 
Based on our results, manufacturers should change their interfaces so integrators can deploy optimal authentication mechanisms.
\end{abstract}

\section{Introduction}
%%%
%%%
%%%
\newcommand{\FAR}{\ensuremath{\textit{FAR}}\xspace}
\newcommand{\FRR}{\ensuremath{\textit{FRR}}\xspace}
\newcommand{\EER}{\ensuremath{\textit{EER}}\xspace}
\newcommand{\ANDintegration}{\ensuremath{
P^{\text{AND}}_{\text{success}}(t)
=
\bigl(1-\FRR(t)\bigr)
\Bigl(
P^{1}_{\text{safe}}
+
P^{1}_{\text{leak}}\bigl(1-\FAR(t)\bigr)
\Bigr)
}\xspace}
\newcommand{\ORintegration}{\ensuremath{
    P^{\text{OR}}_{\text{success}}(t)
=
\bigl(1-\FAR(t)\bigr)
\Bigl(
P^{1}_{\text{safe}}
+
P^{1}_{\text{loss}}\bigl(1-\FRR(t)\bigr)
\Bigr)
}\xspace}
\newcommand{\Psuccess}{\ensuremath{
P_{\text{success}}(t)=
\begin{cases}
P^{1}_{\text{safe}}, 
& t \le a_1, \\[10pt]

(P^{1}_{\text{safe}} + P^{1}_{\text{leak}})
- P^{1}_{\text{leak}}\,\dfrac{a_2 - t}{a_2 - a_1}, 
& a_1 < t \le u_1, \\[14pt]

\dfrac{u_2 - t}{u_2 - u_1}
\left[
(P^{1}_{\text{safe}} + P^{1}_{\text{leak}})
- P^{1}_{\text{leak}}\,\dfrac{a_2 - t}{a_2 - a_1}
\right],
& u_1 < t \le a_2, \\[16pt]

\dfrac{u_2 - t}{u_2 - u_1}
(P^{1}_{\text{safe}} + P^{1}_{\text{leak}}),
& a_2 < t \le u_2, \\[12pt]

0, 
& t > u_2
\end{cases}
}\xspace}
\newcommand{\contprobabilities}{
\begin{align*}
P^{2}_{\text{safe}}(t) &= 1 - \FAR(t) - \FRR(t) + \FAR(t) \cdot \FRR(t),\\
P^{2}_{\text{loss}}(t) &= \FRR(t)\bigl(1 - \FAR(t)\bigr),\\
P^{2}_{\text{leak}}(t) &= \FAR(t)\bigl(1 - \FRR(t)\bigr),\\
P^{2}_{\text{theft}}(t) &= \FAR(t) \cdot \FRR(t).
\end{align*}
}
Biometric credentials are increasingly adopted for access control, as they appear in~81\% of smartphones and~70\% of financial banking services, mainly through face or fingerprint verification~\cite{most_smartphones_biometric,financial_biometric}. They rely on unique physical properties that make them difficult to forge~\cite{biometric_importance}. Biometric credentials are evaluated by a sensor that obtains a~\emph{sample}. It calculates a \emph{similarity score} between that sample and a~\emph{reference} stored in advance, and outputs~\emph{accept} if the score is above a~\emph{match threshold} or \emph{reject} otherwise. This is in contrast to \emph{atomic credentials}, which are either presented correctly or not. Those are evaluated according to a perfect match criterion.

Previous work~$(\S\ref{sec:related})$ analyzed biometric credentials in isolation, primarily focusing on reducing the \emph{False Acceptance Rate} and \emph{False Rejection Rate}~(\FAR, \FRR)~\cite{Anil1997,prabhakar2003,Cappelli2006}. Sarkar et al.~\cite{sarkar2020} suggest to set the threshold heuristically to obtain an Equal Error Rate (\EER, where~$\FAR = \FRR$). But these works do not analyze the integration of biometric credentials into an authentication mechanism. Other works analyze authentication mechanisms as predicates of credential possession~\cite{maram2024,mouallem2024,Eyal2021,Abhishek2013}, but treat all credentials as atomic. 

In this work, we study how to design \emph{optimal mechanisms},  mechanisms with the maximum probability of success given the credentials and their probabilities, including biometric credentials. First, we present a refined definition of the authentication problem with biometric credentials~$(\S\ref{sec:model})$. The system comprises an authentication mechanism and two principals: a~\emph{user} and an~\emph{attacker}. Each principal interacts with the mechanism by presenting credentials.

Biometric sensors aim to assign high scores to samples from legitimate users and low scores to samples from attackers. The scores obtained by comparing legitimate users' samples against the reference follow the \emph{genuine distribution}. The scores obtained by comparing attackers' samples against the reference follow the \emph{impostor distribution}. Consequently, the genuine distribution should be shifted toward higher matching scores compared to the impostor distribution, as shown in Figure~\ref{fig:intro-fig}. These distributions establish two fault probabilities that depend on the threshold. The~\FAR is the probability that the impostor score value is above the threshold~(red volume) and the~\FRR is the probability that the legitimate user's score value is below the threshold~(blue volume).

To identify an optimal mechanism, we derive closed-form expressions for the mechanism's success probability, using first and second order approximations of both principals' biometric score distributions. We start by optimizing a single biometric credential mechanism~$(\S\ref{sec:mechanisms})$. In a single credential mechanism, the mechanism success probability reduces to the probability that the credential is \emph{safe}, i.e., only to the user can use the credential. We identify the optimal threshold of the biometric credential and show that, in general, the optimal threshold differs from the commonly used~\EER threshold.

\begin{figure}[t]
\centering
\begin{tikzpicture}
\begin{axis}[
    scale only axis,
    width=0.8\columnwidth,
    height=0.50\columnwidth,
    xlabel={Matching Score},
    ylabel={Probability Density},
    xmin=0, xmax=1,
    ymin=0,
    xtick=\empty,
    ytick=\empty,
    grid=none,
    legend style={draw=black, fill=white},
    legend pos=north west,
]

% FRR shaded area: user rejected
\addplot[
    fill=blue,
    fill opacity=0.20,
    draw=none,
    forget plot
]
table [x=x, y=FRR] {figs/fig_2continuous/gaussian_far_frr_data.txt}
\closedcycle;

% FAR shaded area: attacker accepted
\addplot[
    fill=red,
    fill opacity=0.20,
    draw=none,
    forget plot
]
table [x=x, y=FAR] {figs/fig_2continuous/gaussian_far_frr_data.txt}
\closedcycle;

% User distribution
\addplot[blue, thick, mark=none]
table [x=x, y=user] {figs/fig_2continuous/gaussian_far_frr_data.txt};
\addlegendentry{User}

% Attacker distribution
\addplot[red, thick, mark=none]
table [x=x, y=attacker] {figs/fig_2continuous/gaussian_far_frr_data.txt};
\addlegendentry{Attacker}

% Threshold line
\addplot[black, dashed, thick, mark=none]
coordinates {(0.5,0) (0.5,4.2)};
\addlegendentry{Threshold}

\end{axis}
\end{tikzpicture}
\caption{User and Attacker score distributions with \FAR and \FRR areas appear in red and blue respectively.}
\label{fig:intro-fig}
\end{figure}

Next, we extend our analysis to a mechanism that combines two credentials~$(\S\ref{sec:two_cred})$. A mechanism that uses two biometric credentials appears in manufacturing industry and in financial services~\cite{multimodal_biometrics}. We evaluate a mechanism comprising two biometric credentials by considering first and second order approximations of the biometric credential distributions. We then use a numerical optimization algorithm to find the optimal thresholds for each biometric credential. The optimal thresholds of the biometric credentials in such a mechanism are not, in general, the same as the optimal thresholds if used individually and differs from the~\EER. This difference arises as we optimize each credential with respect to the entire mechanism.
To demonstrate the potential impact, we present a synthetic experiment of two sampled score distributions and observe a significant reduction in the failure rate, when comparing using the~\EER and the optimal threshold. 

We then turn to analyze mechanisms that use both a biometric credential and an atomic credential. Such mechanisms are prevalent in real-world authentication systems: for example, Google~\cite{google_passkeys_2023} assumes (1)~possession of a trusted device like a phone or computer (Atomic credential) and (2)~a biometric credential. We derive closed-form expressions for the mechanism's success probability. Our symbolic analysis reveals that the relative influence of the atomic and biometric credentials on the mechanism's success probability, depends on their individual properties. In particular, as one credential becomes increasingly reliable, the mechanism success probability depends mainly on that credential. Consequently, changes in the less reliable credential have a smaller effect on the mechanism success probability. 

Motivated by this observation, we identify the conditions under which the integration of a biometric credential with an atomic credential will always enhance the mechanism's success probability~$(\S\ref{sec:proof_success_improvement})$. First, we show that when an atomic credential can only be \emph{safe} or \emph{stolen}, it is not worthwhile to use both the atomic and biometric credentials. This is because a mechanism that uses an atomic credential alone succeeds whenever the credential is safe. However, adding a biometric credential may reject the user even if the atomic credential is safe, without improving the mechanism's success when the atomic credential is stolen. Next, we consider atomic credentials that can also be~\emph{leaked}, where both the user and attacker can use the credential, or \emph{lost}, where neither can use it. We introduce a necessary and sufficient condition we call \emph{asymptotic separability}, of a biometric credential: the~\FAR converges to zero for matching scores above a given threshold, or the~\FRR converges to zero for matching scores below it. This assumption is reasonable since biometric sensors aim to distinguish between legitimate users and attackers. Suppose the atomic credential has non-zero probabilities for the states of \emph{safe} and either \emph{loss} or \emph{leak} and any probability of \emph{stolen}. We prove that every \emph{asymptotically separable} biometric credential that is optimally integrated with such an atomic credential, yields a mechanism whose success probability dominates that of the atomic credential mechanism.

Biometric mechanisms typically use the same fixed hueristic threshold~\cite{biocred_fixed_threshold}. Motivated by this behaviour, we extend our analysis to a mechanism that uses a biometric credential only and allows multiple authentication attempts by each principal~$(\S\ref{sec:multiple_auth_attempts})$. In practice, biometric real-world authentication mechanisms allow principals to retry authentication following a failed attempt. Examples appear in Apple and Samsung devices~\cite{apple_biometric,samsung_biometric}. However, while additional sampling opportunities increase the legitimate user's success probability, they also provide attackers with additional opportunities to authenticate successfully. Under a bounded number of attempts mechanism, the mechanism succeeds if the legitimate user authenticates at least once while the attacker fails across all attempts. We then derive closed-form expressions for the mechanism's success probability and optimize the match thresholds for each number of attempts. Our analysis shows that the optimal threshold increases with the number of attempts. However, system designers may seek to limit the expected number of attempts required for the legitimate user to successfully authenticate. We therefore propose a different approach in which we don't bound the number of attempts, but maximize the success probability subject to a bound on the mean number of attempts. Conversely, we minimize the mean number of attempts subject to a required success probability. These optimization problems reveal the tradeoff between security and usability in the authentication system.

We quantitatively evaluate our framework~$(\S\ref{sec:empirical_evaluation})$ using public biometric data~\cite{LivDet2015}. We employ a standard fingerprint verification algorithm~\cite{sourceafis} to evaluate the resulting matching scores of fingerprints from multiple users. We produce the user distribution by matching fingerprints from the same user, resulting in high matching scores and the attacker distribution by matching fingerprints from different users. Consequently, the resulting attacker distribution concentrates around low score values. As a result, the user and the attacker distributions share a small overlapping area, resulting in a very high success probability. Thus, no significant improvement occurs when comparing the mechanism's success under the~\EER threshold and the optimal threshold. These results suggest that the~\EER threshold is sufficient for this attack scenario.

However, the advantage of threhsold optimization becomes apparent when attackers are presenting forged biometric samples. Thus, we conduct a spoofing attack evaluation using the LivDet database~\cite{LivDet2009,LivDet2011,LivDet2013,LivDet2015}, which includes both genuine fingerprint scans and corresponding spoofing samples. We match each user's spoofing samples against their geniune samples to obtain the imposter score distribution. These spoof samples produce matching scores that closely resemble those of the legitimate user, increasing the biometric credential fault probabilities. Thus, threshold optimization becomes necessarry in a spoofing attack scenario. We show that integrating the biometric credential into a mechanism that uses an atomic credential and optimizing the threshold, reduces the mechanism’s failure probability by~20-72\% compared to its failure probability under the~\EER threshold. Furthermore, the results show a reduction of up to~21\% in the mechanism’s failure probability compared to a mechanism based solely on an atomic credential, demonstrating the benefit of taking advantage of the biometric credential. 

Our previous evaluation demonstrates that optimizing the match threshold significantly improves the mechanism's success probability under spoofing attacks. However, in practice, an attacker may perform multiple authentication attempts using spoofed fingerprints, potentially increasing his probability of successfully authenticating. We therefore extend our evaluation to examine whether threshold optimization remains beneficial under multiple authentication attempts. Using the same user and attacker distributions, we compare the mechanism success probability when the thresholds optimized jointly across attempts with two baselines: reusing the~\EER threshold and reusing the optimal single attempt threshold. Our results show that the mechanism success probability is lower by up to 75\% on both baselines than jointly optimizing the thresholds. Moreover, the optimal thresholds value increase with the number of attempts, as predicted by our theoretical analysis. 

To assess whether commercial biometric sensors support threshold optimization, we review their specifications~$(\S\ref{sec:sensors_review})$. We examine 32 biometric device families from 21 manufacturers, covering different modalities, including fingerprints, face, iris, and palm recognition. Our findings reveal that only two device families permit tuning a numeric match threshold, while the others expose target error rates, discrete security levels, or no threshold control at all. However, setting a threshold is necessary but not sufficient for optimization: the integrator must also know~\FAR and~\FRR functions to choose a threhsold. Applying this criterion reveals that no family publishes these functions. Consequently, no surveyed sensor supports threshold optimization, limiting the deployment of optimal authentication mechanisms.

In summary~$(\S\ref{sec:conclusion})$, our contributions are both theoretical and practical: (1)~Formalization of the success probability function for authentication mechanisms using biometric and atomic credentials; (2)~Symbolic analysis of the optimal threshold for a mechanism using a biometric credential; (3)~Identification of the conditions under which adding a biometric credential improves the mechanism's success probability; (4)~Proof that any biometric credential satisfying these conditions improves the mechanism's success probability; (5)~Derivation of the mechanism's success probability under multiple authentication attempts; (6)~Formulation of constrained optimization problems balancing the mechanism's success probability and the mean number of authentication attempts; (7)~Empirical evaluation of multiple authentication attempts under spoofing attacks using real biometric data; (8)~Practical implication that biometric sensor manufacturers should enable threshold configuration to improve authentication security.

\section{Related Work}\label{sec:related}
\noindent
Authentication is a constant active and prominent research area, with prior work broadly categorized into two main aspects: credentials and mechanisms. Credentials are commonly manifested as passwords~\cite{password1979}, One-Time-Password (OTP)~\cite{otp2017}, physical devices~\cite{FIDO2021} or biometric features such as fingerprints, face recognition or iris~\cite{anil2004}. 

Extensive work has evaluated and improved the accuracy of biometric systems~\cite{Cappelli2006,Maio2002,Anil1997,LivDet2009,LivDet2011,LivDet2013,LivDet2015}. Prabhakar et al.~\cite{prabhakar2003} examine the tradeoff between two key error metrics of a biometric system, namely False Acceptance Rate~(\FAR) and False Rejection Rate~(\FRR). Sarkar et al.~\cite{sarkar2020} define the term Equal Error Rate~(\EER), the threshold where~\FAR and~\FRR are equal. They heuristically suggest it as the threshold for biometric credentials. However, these works do not define a success criterion for authentication mechanisms, nor do they explore how the presence of additional credentials can influence the optimal matching threshold of the biometric system. Hong et al.~\cite{hong1998} integrate fingerprint and face recognition, each one operates independently, using an \emph{AND mechanism}. They design a decision fusion protocol that makes a decision considering the~\FAR error metric only. However, they do not explore how one biometric credential affects the other for optimal usage, nor do they analyze the integration of biometric credentials along with other credential types within an authentication mechanism.

An authentication mechanism is a function that evaluates a set of credentials provided by a principal~\cite{Abhishek2013} and output accept or reject~\cite{MattBishop2004,schneider_auth_people}. Recent work has explored authentication mechanism defined as predicate of credential possession. Maram et al.~\cite{maram2024} study interactive authentication mechanisms, where the user and attacker interact with the authentication mechanism. Mouallem et al.~\cite{mouallem2024} compare mechanisms' success and aim to identify mechanisms with maximal success when composed of a large number of homogeneous credentials. Eyal~\cite{Eyal2021} analyzes mechanisms modeled as Boolean functions of the credentials. However, these works are partially applicable to our setting. While we adopt their mechanism's success definition and the way they define the credential possession states, they model biometric credentials as atomic credentials neglecting the tuneable threshold.

%%%
%%%
%%%
\section{Model}\label{sec:model}
\noindent
We formalize the authentication problem by specifying an execution and the system participants~$(\S\ref{sec:System_participants})$, credential types and their characteritics~$(\S\ref{sec:credentials})$, and the authentication mechanism and its success~$(\S\ref{sec:model_mechanism})$.

\subsection{System Participants}\label{sec:System_participants}
\noindent
The system comprises an authentication mechanism~$M$ and two principals: a user~$U$ and an attacker~$A$. In an execution~$E$, a principal~$P$ \emph{presents} a credential~$c_i$ to the mechanism by calling a random function \emph{present credentials} denoted by~$\textit{present}_P(i)$. We abstract away the concrete implementation of a credential and model only its presentation. The value of the~\emph{presentation function} determines whether a credential is available to a principal. The entry~$\sigma_i^P=1$ means that~$P$ successfully presented credential~$c_i$, and~$\sigma_i^P=0$ means that it is not. At the beginning of an execution~$E$, the system is parameterized such by an \emph{allocation function}. A credential is successfully presented with probability randomized by the allocation function denoted by~$\textit{allocate}_P(i)$.

\begin{definition}(Availability Vector)
    An \emph{availability vector} specifies, for each credential, whether it is successfully presented by a principal~$P$ or not. Denote by $\sigma^U$ \textit{and} $\sigma^A$ the availability vectors of the user and of the attacker, respectively.
    \begin{equation*}
    \sigma_i^P = \textit{present}_P(i)
    \end{equation*}
\end{definition}
The outcome of the presentation function determines whether~$\sigma_i^P=1$ or $\sigma_i^P=0$. A scenario~$\sigma$ is a tuple of availability vectors~$(\sigma^U, \sigma^A)$. 

\subsection{Credentials}\label{sec:credentials}
\noindent
A credential may correspond to secret data, such as a password or one time password~(OTP); a physical object, such as a phone or security key; or a biometric property, such as a fingerprint or face. We model these credentials as either \emph{atomic} or \emph{biometric}. Each credential can be either:~$\textit{Safe}$: Only available to the user, $\textit{Lost}$: Not available to either principal, $\textit{Leaked}$: Available to both principals, or~$\textit{Stolen}$: Only available to the attacker.
\subsubsection*{3.2.1}\label{sec:3.2.1}
Passwords, OTPs and security keys are credential types evaluated according to a perfect match criterion, namely, atomic credentials.
\begin{definition} (Atomic Credential)~For an atomic credential the allocation function~$\textit{allocate}_P(i)$ randomly draws the state probabilities of the credential, namely~$P_{\text{safe}}$, $P_{\text{loss}}$, $P_{\text{leak}}$ and~$P_{\text{stolen}}$.
\end{definition}
When a principal~$P$ presents an atomic credential~$c_i$, the presentation function ouputs~1 with probability~$P_{\text{safe}} + P_{\text{leak}}$ when allocated to the user and 0 with probability~$P_{\text{leak}} + P_{\text{stolen}}$ when allocated to the attacker.
\subsubsection*{3.2.2}\label{sec:3.2.2}
Alternatively, a credential may correspond to a biometric property, such as a fingerprint or face. The biometric sensor evaluates a measurement of the biometric property and produces a score~$s$. The user distribution~$f_U(s)$ denotes the score distribution of samples from the user, whereas the attacker distribution~$f_A(s)$ denotes the score distribution of samples from the attacker.
\begin{definition} (Biometric Credential) For a biometric credential, the user score distribution~$f_U(s)$, the attacker score distribution~$f_A(s)$, and the match threshold~$t$ define the presentation function~$\textit{present}_P(i)$. For each principal~${P\in\{U,A\}}$, the allocation function~$\textit{allocate}_P(i)$ draws a random score~$s$ from the distribution~$f_P(s)$. Thus
    \begin{equation*}
    s = \textit{allocate}_P(i)
    \end{equation*}
    \begin{equation*}
    \textit{present}_P(i) = \begin{cases}
        1 &  s \ge t \\
        0 & \text{otherwise}
    \end{cases}
    \end{equation*}
\end{definition}
We denote by~$\sigma_i$ the state of credential~$c_i$ in scenario~$\sigma$. 
% Each state of the biometric credential corresponds to the outcome of the presentation functions for a user and an attacker. Specifically, the pair $(\textit{present}_U(i), \textit{present}_A(i))$ determines whether the biometric credential is~\emph{Safe}, \emph{Loss}, \emph{Leak} or \emph{Stolen}. Therefore, the state probabilities are
% \begin{equation*}
% P_{\text{safe}} = \Pr[\mathsf{present}_U(i) = 1 \wedge \mathsf{present}_A(i) = 0].
% \end{equation*}
% \begin{equation*}
% P_{\text{loss}} = \Pr[\mathsf{present}_U(i) = 0 \wedge \mathsf{present}_A(i) = 0].
% \end{equation*}
% \begin{equation*}
% P_{\text{leak}} = \Pr[\mathsf{present}_U(i) = 1 \wedge \mathsf{present}_A(i) = 1].
% \end{equation*}
% \begin{equation*}
% P_{\text{stolen}} = \Pr[\mathsf{present}_U(i) = 0 \wedge \mathsf{present}_A(i) = 1].
% \end{equation*}
The user and attacker distributions establish two error metrics:(1)~\FAR, the false acceptance rate, is the probability that the sensor outputs \emph{accept} for an attacker sample. We calculate it by integrating the attacker distribution from the match threshold to the maximum score, thus
\begin{equation}
    \label{eqn:FAR}
\FAR(t) = \int_{t}^{\infty} f_A[s] ds
\end{equation}
(2)~\FRR, the false rejection rate, is the probability that the sensor outputs \emph{reject} for a user sample. We calculate it by integrating the user distribution from the minimum score to the match threshold, thus
\begin{equation}
    \label{eqn:FRR}
\FRR(t) = \int_{-\infty}^{t} f_U[s] ds
\end{equation}
We use~\FAR and~\FRR to find the state probabilities of the biometric credential. The user has access to the biometric credential when the output of the biometric sensor is \emph{accept}, which occurs with probability~$1-\FRR(t)$. The attacker has access to the biometric credential when the output of the biometric sensor is \emph{accept}, which occurs with probability~$\FAR(t)$. Thus, the probabilities of the four states are
\begin{equation*}
    P_{\text{loss}} = \FRR(t) \cdot (1 - \FAR(t)),
\end{equation*}
\begin{equation*}
P_{\text{leak}} = \FAR(t) \cdot (1 - \FRR(t)),
\end{equation*}
\begin{equation*}
P_{\text{stolen}} = \FRR(t) \cdot \FAR(t), \textit{ and }
\end{equation*}
\begin{equation}\label{eqn:safe_prob}
P_{{\text{safe}}} =1-\mathrm{P_{\text{loss}}}-\mathrm{P_{\text{leak}}}-\mathrm{P_{\text{stolen}}} = (1 - \FAR(t)) \cdot (1 - \FRR(t)).
\end{equation}

% For example, if~$\sigma_i = \text{leak}$ then the credential is available to both players, namely~$\sigma_i^U = \text{True}$ and~$\sigma_i^A = \text{True}$. If~$c_i$ is an atomic credential, we can deduce that~$x_{c_i^{\mathrm{atm}}}^{U}=1$ $\ \text{and}\ x_{c_i^{\mathrm{atm}}}^{A}=1$. In the case where~$c_i$ is a biometric credential, this implies~$s_{i}^U \ge t \ \text{and}\ s_{i}^A \ge t$. The probabilities that a user or an attacker satisfies an atomic credential can be characterized using password statistics from company databases~\cite{password_leak}. In contrast, the probabilities that a user or an attacker satisfies a biometric credential are characterized by the score distributions of the biometric credential.
\subsection{Mechanism}\label{sec:model_mechanism}
\noindent
An authentication mechanism~$M$ receives the~\emph{evaluation} of the presentation function~$\textit{present}_P(i)$ for each principal~$P$ and outputs a \emph{decision}.
\begin{definition} (Authentication Mechanism) An authentication mechanism $M$ is a predicate that takes as input:
\begin{itemize}
     \item~$\sigma^P$ the availability vector of a principal~$P$ in scenario~$\sigma$ (perhaps an empty vector, represented by~$\varnothing$).
\end{itemize}
and outputs either~1, meaning that the mechanism accepts the principal or~0, meaning that the mechanism rejects the principal.
\end{definition}

We are now ready to define the mechanism's success.
The mechanism is successful if it grants access to the user only. 
\begin{definition} (Mechanism Success) A mechanism~$M$ is \emph{successful} in a scenario~$\sigma$ if for a principal~$P$ and its availability vector~$\sigma^P$, the authentication mechanism decision~$\mathrm{decide}_M(\sigma^P)$ equals~1 for the user only over an execution~$E$, that is
\begin{equation*}
    \mathrm{decide}_M(\sigma^U)=1 \wedge \mathrm{decide}_M(\sigma^A)=0.
\end{equation*}
Otherwise, the mechanism is \emph{unsuccessful}.
\end{definition} 
The mechanism's \emph{success probability} is thus
\begin{equation*}
P_{\text{Success}}(M) =
\Pr\big[\mathrm{decide}_M(\sigma^U)=1 \wedge \mathrm{decide}_M(\sigma^A)=0\big].
\end{equation*}
For example, in an \emph{AND mechanism}~\cite{Eyal2021}, the mechanism grants access to a principal if and only if all credentials are available to the principal; thus,~$\mathrm{decide}_M(\sigma^P)=1$ if ~$\sigma^P_i$ equals~1 for every credential~$c_i$. In an \emph{OR mechanism}, the mechanism grants access to a principal if at least one credential is available to the principal; thus,~$\mathrm{decide}_M(\sigma^P)=1$ if ~$\sigma^P_i$ equals~1 for at least one credential~$c_i$.

\section{Real Data and Approximations}
\label{sec:real_data}
\noindent
We evaluate our model using~\textit{LivDet2015}, a real-world dataset~\cite{LivDet2015,LivDet2013,LivDet2011,LivDet2009}.
It contains multiple fingerprint scans of each user, acquired using different types of sensors, as well as spoof samples produced from each user's fingerprints.
We divide each user's scans into reference scans and candidate scans. 
To score the fingerprint scan images, we use a standard fingerprint verification algorithm, namely~\textit{SourceAFIS}~\cite{sourceafis}. 
We compute a similarity score for each candidate scan by comparing it against every reference image, and retaining the maximum value as the candidate's score~\cite{hong1998,Cappelli2006,max_score2009}.
For each of two attack scenarios, we produce histograms of genuine and impostor scores~\cite{Anil1997,Cappelli2006,Slob2005}, to obtain the empirical user and attacker score distributions.

In the first scenario, an attacker authenticates using their own fingerprints. 
The impostor scores are all the users' scans evaluated against the reference of the other users in the database.
With an~\EER of~3.8\%, the algorithm’s performance is representative of state-of-the-art fingerprint verification~\cite{Maio2002,Cappelli2006} and the best results achieved in the~\textit{LivDet2015} competition~\cite{LivDet2015}.
In the second scenario, an attacker creates fake fingerprint samples using spoofing materials such as Gelatine, Woodglue, Latex and Ecoflex. We match each spoof sample against the spoofed user's reference images to obtain the impostor score distribution.

To analyze different mechanisms, we approximate the empirical score distributions. 
We evaluate the accuracy by fitting uniform, quadratic, and normal distributions to the genuine and impostor score histograms in both attack scenarios.
For each fitted distribution, we calculate the Kolmogorov-Smirnov~(KS) statistic~\cite{Massey1951} against the~\FRR and~\FAR curves.
Figure~\ref{fig:approximations} presents the resulting~\FAR and~\FRR curves alongside the empirical values for both scenarios.
The KS statistics range from~$0.04$ to~$0.07$ for the genuine scores and are~$0.17$ for the impostor scores, providing empirical support for using our theoretical models as approximations of the measured score distributions.

\begin{figure*}[t]
\centering
\pgfplotsset{
    impostor curve/.style={line width=0.6pt, dashed},
    genuine curve/.style={line width=0.6pt, solid},
    uniform model/.style={purple},
    quadratic model/.style={orange},
    normal model/.style={cyan},
    ksaxis/.style={
        scale only axis,
        width=0.86\linewidth,
        height=0.40\linewidth,
        xlabel={Threshold (\%)},
        ylabel={Error Rate},
        xmin=0, xmax=100,
        ymin=0, ymax=1,
        grid=major
    }
}

\begin{subfigure}[t]{0.48\textwidth}
\centering
\begin{tikzpicture}
\begin{axis}[
    ksaxis,
    legend to name=kslegend,
    legend columns=6,
    legend style={draw=black, fill=white, font=\scriptsize}
]
\addplot[only marks, red, mark=*, mark size=1.1pt]
table[x=T, y=FAR] {figs/fig_different_users_sourceafsi/sourceafis_far_frr_data.txt};
\addlegendentry{Empirical \FAR}

\addplot[only marks, blue, mark=*, mark size=1.1pt]
table[x=T, y=FRR] {figs/fig_different_users_sourceafsi/sourceafis_far_frr_data.txt};
\addlegendentry{Empirical \FRR}

\addplot[impostor curve, uniform model]
table[x=T, y=FAR_uniform] {figs/fig_different_users_sourceafsi/sourceafis_far_frr_data.txt};
\addlegendentry{Uniform \FAR}

\addplot[genuine curve, uniform model]
table[x=T, y=FRR_uniform] {figs/fig_different_users_sourceafsi/sourceafis_far_frr_data.txt};
\addlegendentry{Uniform \FRR}

\addplot[impostor curve, quadratic model]
table[x=T, y=FAR_quadratic] {figs/fig_different_users_sourceafsi/sourceafis_far_frr_data.txt};
\addlegendentry{Quadratic \FAR}

\addplot[genuine curve, quadratic model]
table[x=T, y=FRR_quadratic] {figs/fig_different_users_sourceafsi/sourceafis_far_frr_data.txt};
\addlegendentry{Quadratic \FRR}

% \addplot[impostor curve, normal model]
% table[x=T, y=FAR_normal] {figs/fig_different_users_sourceafsi/sourceafis_far_frr_data.txt};
% \addlegendentry{Normal \FAR}

% \addplot[genuine curve, normal model]
% table[x=T, y=FRR_normal] {figs/fig_different_users_sourceafsi/sourceafis_far_frr_data.txt};
% \addlegendentry{Normal \FRR}
\end{axis}
\end{tikzpicture}
\caption{Different users.}
\label{fig:ks_different_users}
\end{subfigure}
\hfill
\begin{subfigure}[t]{0.48\textwidth}
\centering
\begin{tikzpicture}
\begin{axis}[ksaxis]
\addplot[only marks, red, mark=*, mark size=1.1pt]
table[x=T, y=FAR] {figs/fig_spoofed_users/livdet_sourceafis_far_frr_best_of_probes_data.txt};

\addplot[only marks, blue, mark=*, mark size=1.1pt]
table[x=T, y=FRR] {figs/fig_spoofed_users/livdet_sourceafis_far_frr_best_of_probes_data.txt};

\addplot[impostor curve, uniform model]
table[x=T, y=FAR_uniform] {figs/fig_spoofed_users/livdet_sourceafis_far_frr_best_of_probes_data.txt};

\addplot[genuine curve, uniform model]
table[x=T, y=FRR_uniform] {figs/fig_spoofed_users/livdet_sourceafis_far_frr_best_of_probes_data.txt};

\addplot[impostor curve, quadratic model]
table[x=T, y=FAR_quadratic] {figs/fig_spoofed_users/livdet_sourceafis_far_frr_best_of_probes_data.txt};

\addplot[genuine curve, quadratic model]
table[x=T, y=FRR_quadratic] {figs/fig_spoofed_users/livdet_sourceafis_far_frr_best_of_probes_data.txt};

% \addplot[impostor curve, normal model]
% table[x=T, y=FAR_normal] {figs/fig_spoofed_users/livdet_sourceafis_far_frr_best_of_probes_data.txt};

% \addplot[genuine curve, normal model]
% table[x=T, y=FRR_normal] {figs/fig_spoofed_users/livdet_sourceafis_far_frr_best_of_probes_data.txt};
\end{axis}
\end{tikzpicture}
\caption{Spoofing attack.}
\label{fig:ks_spoofing}
\end{subfigure}

\medskip
\pgfplotslegendfromname{kslegend}

\caption{Empirical~\FAR and~\FRR as functions of the threshold, compared
with the uniform, quadratic and normal distribution functions obtained
from the KS fits, under both attack scenarios.}
\label{fig:approximations}
\end{figure*}

\section{Single Biometric Credential}\label{sec:mechanisms}
\noindent
We explore optimal mechanism success probability with a single biometric credential.
We start with first~$(\S\ref{sec:uniform_biocred})$ and second~$(\S\ref{sec:quadratic_biocred})$ order approximations, and then evaluate the mechanism on public biometric datasets~$(\S\ref{sec:eval_single_cred})$.

\begin{figure*}[t]
\centering

\begin{subfigure}[t]{0.30\textwidth}
\centering
\begin{tikzpicture}
\begin{axis}[
    scale only axis,
    width=0.8\linewidth,
    height=0.5\linewidth,
    xlabel={Matching Score},
    ylabel={Probability Density},
    xmin=0, xmax=1,
    ymin=0, ymax=4.8,
    grid=major,
    xtick={0,0.15,0.40,0.60,0.85},
    xticklabels={$0$,$a_1$,$a_2$,$u_1$,$u_2$},
    ytick=\empty,
    legend to name=uniformthreecaseslegend,
    legend columns=2,
    legend style={draw=black, fill=white}
]
\addplot+[blue, thick, mark=none, const plot]
table [x=s, y=PU_case1] {figs/fig_uniform/uniform_three_cases.txt};
\addlegendentry{User}

\addplot+[red, thick, mark=none, const plot]
table [x=s, y=PA_case1] {figs/fig_uniform/uniform_three_cases.txt};
\addlegendentry{Attacker}
\end{axis}
\end{tikzpicture}
\caption{Separated uniform PDFs.}
\label{fig:uniform_case1_pdf}
\end{subfigure}
\begin{subfigure}[t]{0.30\textwidth}
\centering
\begin{tikzpicture}
\begin{axis}[
    scale only axis,
    width=0.8\linewidth,
    height=0.5\linewidth,
    xlabel={Matching Score},
    ylabel={},
    xmin=0, xmax=1,
    ymin=0, ymax=4.8,
    grid=major,
    xtick={0,0.20,0.40,0.65,0.85},
    xticklabels={$0$,$u_1$,$a_1$,$a_2$,$u_2$},
    ytick=\empty
]
\addplot+[blue, thick, mark=none, const plot]
table [x=s, y=PU_case2] {figs/fig_uniform/uniform_three_cases.txt};

\addplot+[red, thick, mark=none, const plot]
table [x=s, y=PA_case2] {figs/fig_uniform/uniform_three_cases.txt};
\end{axis}
\end{tikzpicture}
\caption{The attacker range lies within the user range.}
\label{fig:uniform_case2_pdf}
\end{subfigure}
\begin{subfigure}[t]{0.30\textwidth}
\centering
\begin{tikzpicture}
\begin{axis}[
    scale only axis,
    width=0.8\linewidth,
    height=0.5\linewidth,
    xlabel={Matching Score},
    ylabel={},
    xmin=0, xmax=1,
    ymin=0, ymax=4.8,
    grid=major,
    xtick={0,0.15,0.40,0.65,0.85},
    xticklabels={$0$,$a_1$,$u_1$,$a_2$,$u_2$},
    ytick=\empty
]
\addplot+[blue, thick, mark=none, const plot]
table [x=s, y=PU_case3] {figs/fig_uniform/uniform_three_cases.txt};

\addplot+[red, thick, mark=none, const plot]
table [x=s, y=PA_case3] {figs/fig_uniform/uniform_three_cases.txt};
\end{axis}
\end{tikzpicture}
\caption{Partial overlap of the distributions.}
\label{fig:uniform_case3_pdf}
\end{subfigure}

\ref*{uniformthreecaseslegend}

\caption{Quantitative illustration of three distinct cases of uniform attacker and user score distributions. In case (1), the distributions are separated. Case (2) illustrates the case where an attacker is fully contained in the user distribution, and (3) illustrates two overlapping configurations.}
\label{fig:uniform_three_cases}
\end{figure*}

\subsection{Uniformly Distributed Biometric Credentials}\label{sec:uniform_biocred}
\noindent
We consider uniform distribution as a first-order approximation of the biometric credential distributions.
Denote by~$u_1 \textit{ and } u_2$ the user's PDF bounds and by~$a_1 \textit{ and } a_2$  the attacker's PDF bounds as shown in Figure~\ref{fig:uniform_three_cases}. The probability distribution functions for the user and the attacker, $f_U$ and $f_A$ are thus:
\begin{align*}
&f_U(s) = \frac{1}{u_2-u_1} \quad \text{for } u_1 \leq s \leq u_2 \text{ and}\\
&f_A(s) = \frac{1}{a_2-a_1} \quad \text{for } a_1 \leq s \leq a_2.
\end{align*}

The error metrics~\FAR and~\FRR are both functions of the match threshold~$t$ and given by Equations~\ref{eqn:FAR} and~\ref{eqn:FRR}. Thus:

\[
\FAR(t) = \int_{t}^{a_2} \frac{1}{a_2-a_1} ds =
\left\{
\begin{array}{ll}
1 & \text{if } t \le a_1 \\
\frac{a_2 - t}{a_2 - a_1} & \text{if } a_1 < t \le a_2 \\
0 & \text{if } t > a_2
\end{array}
\right.
\]
\[\textit{and}\]
\[
\FRR(t) = \int_{u_1}^{t} \frac{1}{u_2-u_1} ds =
\left\{
\begin{array}{ll}
0 & \text{if } t \le u_1 \\
\frac{t - u_1}{u_2 - u_1} & \text{if } u_1 < t \le u_2 \\
1 & \text{if } t > u_2
\end{array}
\right.
\]

For a biometric credential to be useful, users are expected to produce higher matching scores than attackers. Otherwise, the attacker would be more likely to produce scores that lead to successful authentication, rendering the biometric credential ineffective. Thus, we assume that~${a_2 < u_2}$. There are three distinct cases: 

(1)~$a_2 < u_1$~(Figure~\ref{fig:uniform_case1_pdf}): In this case, the attacker and user distributions are \emph{separated}. For a mechanism with a single credential~$M^{\textit{single}}$, the mechanism's success probability is simply the probability that the credential is safe
\begin{equation}
    \label{eqn:safesuccess}
    P_{\text{success}}(M^{\text{single}}) = P_{{\text{safe}}}
\end{equation} 
Any threshold~$a_2 \le t \le u_1$ yields zero~\FAR and~\FRR, and therefore achieves a mechanism success probability of~$1$.

(2)~$u_1 < a_1 < a_2 < u_2$~(Figure~\ref{fig:uniform_case2_pdf}): The attacker distribution bounds lie within the range of the user distribution bounds. 
The success probability of the biometric credential is given by Equation~\ref{eqn:safe_prob}.
The~\FAR is one for any~$t \le a_1$ and zero for any~$t \ge a_2$.
Thus:
\begin{equation*}
P_{\text{success}}(t) = 
\begin{cases}
    0, & t \le a_1 \\
    (1 - \frac{a_2-t}{a_2-a_1})\cdot(1-\frac{t-u_1}{u_2-u_1}), & a_1 < t \le a_2\\
    1 - \frac{t-u_1}{u_2-u_1}, & a_2 < t \le u_2 \\
    0, & t > u_2
\end{cases}
\end{equation*}

(3)~$a_1 < u_1 < a_2 < u_2$~(Figure~\ref{fig:uniform_case3_pdf}): The user and attacker distributions partially overlap.

We maintain the same approach as the case~(2) to calculate the probability of success. The~\FAR is one for any~$t \le a_1$ and zero for any~$t \ge a_2$, and the~\FRR is zero for any~$t \le u_1$ and one for any~$t \ge u_2$. Thus, in this case, we obtain five regions of the success probability function:
\begin{equation*}
P_{\text{success}}(t) = 
\begin{cases}
    0, & t \le a_1 \\
    1 - \frac{a_2-t}{a_2-a_1}, & a_1 < t \le u_1 \\
    (1 - \frac{a_2-t}{a_2-a_1})\cdot(1-\frac{t-u_1}{u_2-u_1}), & u_1 < t \le a_2\\
    1 - \frac{t-u_1}{u_2-u_1}, & a_2 < t \le u_2 \\
    0, & t > u_2.
\end{cases}
\end{equation*}
In case~(3) the success probability is zero in the first interval, linearly increasing in the second, quadratic in the third, linearly decreasing in the fourth, and zero in the fifth. Case~(2) has the same structure without the linearly increasing second interval.
The function is continuous, so the global maximum lies in the third interval,~$u_1 \le t \le a_2$, where both~\FAR and~\FRR are strictly between zero and one.
Differentiating there gives an interior critical point at $t_{\text{opt}} = \frac{a_1+u_2}{2}$, clipped to the interval boundaries, with
\begin{equation*}
P_{\text{success}}(t_{\text{opt}}) = \frac{(a_1 - u_2)^2}{4 (a_1 - a_2) (u_1 - u_2)}.
\end{equation*}
Figure~\ref{fig:uniform_success_1_cred} illustrates this analysis.
The optimum shows that a wider attacker distribution, meaning a larger~$(a_2 - a_1)$, reduces the mechanism's success probability.

The~\EER threshold is
\begin{equation*}
t_{\EER} = \frac{a_1 u_1 - a_2 u_2}{a_1 - a_2 + u_1 - u_2},
\end{equation*}
which in general differs from~$t_{\text{opt}}$ in both cases.
The two coincide only when the distributions are symmetric, when~${a_2 - a_1 = u_2 - u_1}$, as Figure~\ref{fig:FAR_FRR_uniform_1_cred} illustrates.
The full interval analysis and the boundary cases are in the full version of this report~\cite{biometric_full_version}.

\begin{figure}[t]
\centering
\begin{tikzpicture}
\begin{axis}[
    scale only axis,
    width=0.8\columnwidth,
    height=0.42\columnwidth,
    xlabel={$t$},
    ylabel={$P_{\mathrm{success}}(t)$},
    grid=none,
    xmin=0, xmax=0.75,
    ymin=0, ymax=0.8,
    legend style={
        at={(0.98,0.98)},
        anchor=north east,
        fill=white,
        fill opacity=0.85,
        font=\scriptsize,
        text opacity=1,
        draw=black
    },
    xtick={0,0.05,0.30,0.375,0.42,0.70},
    xticklabels={$0$,$a_1$,$u_1$,$t_{\mathrm{opt}}$,$a_2$,$u_2$},
    ytick={0,0.7137},
    yticklabels={
        $0$
    }
]

\addplot[thick, myPurple, mark=none]
table [x=s, y=safe] {figs/fig_uniform/success_uniform_data.txt};

\addlegendentry{$P_{\mathrm{success}}(t)$}

\end{axis}
\end{tikzpicture}
\caption{Success probability as a function of the threshold.}
\label{fig:uniform_success_1_cred}
\end{figure}

\begin{figure}[t]
\centering
\begin{tikzpicture}
\begin{axis}[
    scale only axis,
    width=0.7\columnwidth,
    height=0.5\columnwidth,
    xlabel={$\FRR$},
    ylabel={$\FAR$},
    grid=major,
    xmin=-0.05, xmax=1.05,
    ymin=-0.05, ymax=1.05,
    xtick={0,0.2,0.4,0.6,0.8,1},
    ytick={0,0.5,1},
    legend style={
        at={(0.98,0.98)},
        anchor=north east,
        fill=white,
        fill opacity=0.85,
        font=\scriptsize,
        text opacity=1,
        draw=black
    }
]

\addplot[thick, myPurple, mark=none]
table [x=FRR, y=FAR] {figs/fig_uniform/success_uniform_data.txt};
\addlegendentry{Asymmetric Curve}

\addplot[thick, dashed, black, mark=none]
table [x=FRR_sym, y=FAR_sym] {figs/fig_uniform/success_uniform_data.txt};
\addlegendentry{Symmetric Curve}

\addplot[only marks, mark=*, mark size=1.3pt, blue]
table [x=FRR_opt, y=FAR_opt] {figs/fig_uniform/success_uniform_data.txt};
\addlegendentry{Optimal Point}

\addplot[only marks, mark=*, mark size=1.3pt, red]
table [x=FRR_eer, y=FAR_eer] {figs/fig_uniform/success_uniform_data.txt};
\addlegendentry{\EER Point}

\addplot[only marks, mark=*, mark size=1.3pt, blue]
table [x=FRR_opt_sym, y=FAR_opt_sym] {figs/fig_uniform/success_uniform_data.txt};


\addplot[only marks, mark=*, mark size=1.3pt, red]
table [x=FRR_eer_sym, y=FAR_eer_sym] {figs/fig_uniform/success_uniform_data.txt};

\end{axis}
\end{tikzpicture}
\caption{\FAR vs. \FRR operating curves for the asymmetric and symmetric distributions in~(3) case, with optimal and~\EER points.}
\label{fig:FAR_FRR_uniform_1_cred}
\end{figure}

\subsection{Quadratic Distributed Credential}\label{sec:quadratic_biocred}
\noindent
\begin{figure*}[t]
\centering

\begin{subfigure}[t]{0.30\textwidth}
\centering
\begin{tikzpicture}
\begin{axis}[
    scale only axis,
    width=0.8\linewidth,
    height=0.5\linewidth,
    xlabel={Matching Score},
    ylabel={Probability Density},
    xmin=0, xmax=1,
    ymin=0, ymax=7,
    grid=major,
    xtick={0,0.15,0.40,0.60,0.85},
    xticklabels={$0$,$\beta_1$,$\beta_2$,$\alpha_1$,$\alpha_2$},
    ytick=\empty,
    legend to name=quadraticthreecaseslegend,
    legend columns=2,
    legend style={draw=black, fill=white}
]
\addplot+[blue, thick, mark=none]
table [x=s, y=PU_case1] {figs/fig_parabola/quadratic_three_cases.txt};
\addlegendentry{User}

\addplot+[red, thick, mark=none]
table [x=s, y=PA_case1] {figs/fig_parabola/quadratic_three_cases.txt};
\addlegendentry{Attacker}
\end{axis}
\end{tikzpicture}
\caption{Separated quadratic PDFs.}
\label{fig:quadratic_case1_pdf}
\end{subfigure}
\begin{subfigure}[t]{0.30\textwidth}
\centering
\begin{tikzpicture}
\begin{axis}[
    scale only axis,
    width=0.8\linewidth,
    height=0.5\linewidth,
    xlabel={Matching Score},
    ylabel={},
    xmin=0, xmax=1,
    ymin=0, ymax=7,
    grid=major,
    xtick={0,0.20,0.40,0.65,0.85},
    xticklabels={$0$,$\alpha_1$,$\beta_1$,$\beta_2$,$\alpha_2$},
    ytick=\empty
]
\addplot+[blue, thick, mark=none]
table [x=s, y=PU_case2] {figs/fig_parabola/quadratic_three_cases.txt};

\addplot+[red, thick, mark=none]
table [x=s, y=PA_case2] {figs/fig_parabola/quadratic_three_cases.txt};
\end{axis}
\end{tikzpicture}
\caption{The attacker roots lie within the user roots.}
\label{fig:quadratic_case2_pdf}
\end{subfigure}
\begin{subfigure}[t]{0.30\textwidth}
\centering
\begin{tikzpicture}
\begin{axis}[
    scale only axis,
    width=0.8\linewidth,
    height=0.5\linewidth,
    xlabel={Matching Score},
    ylabel={},
    xmin=0, xmax=1,
    ymin=0, ymax=7,
    grid=major,
    xtick={0,0.15,0.40,0.65,0.85},
    xticklabels={$0$,$\beta_1$,$\alpha_1$,$\beta_2$,$\alpha_2$},
    ytick=\empty
]
\addplot+[blue, thick, mark=none]
table [x=s, y=PU_case3] {figs/fig_parabola/quadratic_three_cases.txt};

\addplot+[red, thick, mark=none]
table [x=s, y=PA_case3] {figs/fig_parabola/quadratic_three_cases.txt};
\end{axis}
\end{tikzpicture}
\caption{Partial overlap of the distributions.}
\label{fig:quadratic_case3_pdf}
\end{subfigure}

\ref*{quadraticthreecaseslegend}

\caption{Illustration of three distinct cases of quadratic attacker and user score distributions. In case (1), the distributions are separated. Case (2) illustrates the case where the attacker distribution is fully contained within the user distribution, and case (3) illustrates partial overlap between the two distributions.}
\label{fig:quadratic_three_cases}
\end{figure*}
Our goal is to model credentials whose distributions closely resemble those observed in real measurement data. To this end, we analyze a quadratic probability distribution function, which provides a second order approximation of the biometric credential that can also be seen in real fingerprint measurement data~$(\S\ref{sec:real_data})$. This choice is also consistent with the findings of Slobodan et al.~\cite{Slob2005} on Eigenfinger and Eigenpalm feature distributions. Accordingly, we define two probability density functions (PDFs). Let~$f_U$ denote the user distribution with roots~$\alpha_{1}, \alpha_{2}$ and let~$f_A$ denote the attacker distribution with roots~$\beta_{1}, \beta_{2}$. Denote by $a, b, c$ and $d, e, f$ the quadratic coefficients of the user and attacker distributions, respectively, hence:
\[
f_U(s) = 
\left\{
\begin{array}{ll}
    a \cdot s^2+b \cdot s+c & \alpha_1 \leq s \leq \alpha_2 \\
    0 & s > \alpha_2 \lor s < \alpha_1
    \end{array}
\right.
\]
\[
f_A(s) =
\left\{
\begin{array}{ll}
    d \cdot s^2+e \cdot s+f & \beta_1 \leq s \leq \beta_2 \\
    0 & s > \beta_2 \lor s < \beta_1
    \end{array}
\right.
\]
To reduce the degrees of freedom, we express the coefficients in terms of the roots and set the integral of each PDF over its roots interval to~1. This normalization condition ensures that the resulting error metrics~\FAR and~\FRR are bounded by~1, and determines the coefficients fully. The derivation of the coefficients is in the full version of this report~\cite{biometric_full_version}. 
Therefore, we can write the quadratic PDFs as
\begin{flalign*}
&f_U(s) =
\left\{
\begin{array}{@{}l@{\qquad}l@{}}
     \dfrac{6(s-\alpha_1)(s-\alpha_2)}{(\alpha_1-\alpha_2)^3} & \alpha_1 \leq s \leq \alpha_2 \\
     0, & \text{otherwise}
\end{array}
\right.
&&
\end{flalign*}
\begin{flalign*}
&f_A(s) =
\left\{
\begin{array}{@{}l@{\qquad}l@{}}
     \dfrac{6(s-\beta_1)(s-\beta_2)}{(\beta_1-\beta_2)^3} & \beta_1 \leq s \leq \beta_2 \\
     0, & \text{otherwise}
\end{array}
\right.
&&
\end{flalign*}
These distributions define the error metrics~\FAR and~\FRR as functions of the match threshold~$t$. Denote
\begin{equation*}
A(t)=
\frac{
\begin{aligned}
&2(\beta_2^3-t^3)-3(\beta_1+\beta_2)(\beta_2^2-t^2)\\
&\quad+6\beta_1\beta_2(\beta_2-t)
\end{aligned}
}{
(\beta_1-\beta_2)^3
},
\end{equation*}
\begin{equation*}
U(t)=
\frac{
\begin{aligned}
&2(t^3-\alpha_1^3)-3(\alpha_1+\alpha_2)(t^2-\alpha_1^2)\\
&\quad+6\alpha_1\alpha_2(t-\alpha_1)
\end{aligned}
}{
(\alpha_1-\alpha_2)^3
}.
\end{equation*}

\begin{equation*}
\FAR(t)=
\begin{cases}
1, & t \le \beta_1,\\
A(t), & \beta_1 < t \le \beta_2,\\
0, & t > \beta_2,
\end{cases}
\end{equation*}
\[\textit{and}\]
\begin{equation*}
\FRR(t)=
\begin{cases}
0, & t \le \alpha_1,\\
U(t), & \alpha_1 < t \le \alpha_2,\\
1, & t > \alpha_2.
\end{cases}
\end{equation*}

Similarly to the analysis of the uniform distribution, the x axis represents the matching score $s$. Genuine users are expected to produce higher matching scores than attackers, thus we analyze three different cases:

(1)~$\beta_{2} < \alpha_{1}$~(Figure~\ref{fig:quadratic_case1_pdf}): The two distributions are separated. Any threshold $t$ chosen within the range~$\beta_{2} \le t \le \alpha_{1}$ yields zero~\FAR and~\FRR, and therefore achieves a mechanism success probability of~$1$.

(2)~$\alpha_{1} < \beta_{1} < \beta_{2} < \alpha_{2}$~(Figure~\ref{fig:quadratic_case2_pdf}): The attacker distribution bounds lie within the range of the user distribution bounds. The mechanism’s probability of success reduces to the safe probability of the biometric credential, which is given by Equation~\ref{eqn:safe_prob}.
The~\FAR is one for any~$t \le \beta_1$ and zero for any~$t \ge \beta_2$. Thus
\begin{equation*}
\begin{aligned}
G(t)=
-&\frac{
(t-\alpha_2)^2(t-\beta_1)^2
}{
(\alpha_1-\alpha_2)^3(\beta_1-\beta_2)^3
} \\
&\cdot
(2t-3\alpha_1+\alpha_2)
(2t+\beta_1-3\beta_2).
\end{aligned}
\end{equation*}
\begin{multline*}
P_{\mathrm{success}}(t)=\\
\begin{cases}
0,
& t<\beta_1 \lor t>\alpha_2,\\
G(t),
& \beta_1 \le t \le \beta_2,\\
-\dfrac{(t-\alpha_2)^2(2t-3\alpha_1+\alpha_2)}
{(\alpha_1-\alpha_2)^3},
& \beta_2 \le t \le \alpha_2.
\end{cases}
\end{multline*}

(3)~$\beta_{1} < \alpha_{1} < \beta_{2} < \alpha_{2}$~(Figure~\ref{fig:quadratic_case3_pdf}): The user and attacker distributions partially overlap. We maintain the same approach as in case~(2) to calculate the probability of success. The~\FAR is one for any~$t \le \beta_1$ and zero for any~$t \ge \beta_2$, and the~\FRR is zero for any~$t \le \alpha_1$ and one for any~$t \ge \alpha_2$. Thus, in this case, we obtain four regions of the success probability function
\begin{multline*}
P_{\mathrm{success}}(t)=\\
\begin{cases}
0,
& t<\beta_1 \lor t>\alpha_2,\\
\dfrac{(t-\beta_1)^2(2t+\beta_1-3\beta_2)}
{(\beta_1-\beta_2)^3},
& \beta_1 \le t \le \alpha_1,\\
G(t),
& \alpha_1 \le t \le \beta_2,\\
-\dfrac{(t-\alpha_2)^2(2t-3\alpha_1+\alpha_2)}
{(\alpha_1-\alpha_2)^3},
& \beta_2 \le t \le \alpha_2.
\end{cases}
\end{multline*}
The second and third intervals in case~(2) are identical to the third and fourth intervals in case~(3) and therefore, we analyze them together. 
We differentiate~$P_{\text{success}}(t)$ and compute its critical points in every interval. By comparing the stationary values obtained in all valid regions, we identify the global maximum, which corresponds to the optimal threshold~$t_{\text{opt}}$.
We differentiate~$P_{\text{success}}(t)$ in every interval. In the second and fourth intervals the only critical points lie at~$t=\beta_1$ and~$t=\alpha_2$, where the derivative does not change sign in the required direction, and therefore these intervals contain no maximum. We turn to the third interval, where the derivative is a cubic polynomial multiplied by two linear factors that vanish outside the interval. Thus, the critical points depend solely on the cubic term, which reveals exactly one real root with the remaining two roots being complex. By examining the sign of the derivative on both sides of this root, we verify that it constitutes the global maximum~$t_{\text{opt}}$. The closed form cannot be simplified beyond a complex expression~(see Appendix~\ref{app:quad_opt_thr}), and the full derivative analysis is in the full version of this report~\cite{biometric_full_version}.

% \begin{figure*}[t]
% \centering

% \begin{subfigure}[t]{0.31\textwidth}
% \centering
% \begin{tikzpicture}
% \begin{axis}[
%     scale only axis,
%     width=0.8\linewidth,
%     height=0.45\linewidth,
%     xlabel={$s$},
%     ylabel={$P(s)$},
%     xmin=1.5, xmax=4.5,
%     ymin=-0.1, ymax=1.65,
%     xtick={2,3,3.5,4},
%     xticklabels={$\beta_1$,$\alpha_1$,$\beta_2$,$\alpha_2$},
%     ytick=\empty,
%     legend to name=distlegend,
%     legend columns=2,
%     legend style={draw=black, fill=white}
% ]
% \addplot[blue, thick, mark=none]
% table [x=s, y=PU] {figs/fig_parabola/parabola_case_a.txt};
% \addlegendentry{User}

% \addplot[red, thick, mark=none]
% table [x=s, y=PA] {figs/fig_parabola/parabola_case_a.txt};
% \addlegendentry{Attacker}
% \end{axis}
% \end{tikzpicture}
% \caption{Attacker with low variance.}
% \end{subfigure}
% \hfill
% \begin{subfigure}[t]{0.31\textwidth}
% \centering
% \begin{tikzpicture}
% \begin{axis}[
%     scale only axis,
%     width=0.8\linewidth,
%     height=0.45\linewidth,
%     xlabel={$s$},
%     ylabel={$P(s)$},
%     xmin=1, xmax=4.5,
%     ymin=-0.1, ymax=1.65,
%     xtick={1.2,3,3.5,4},
%     xticklabels={$\beta_1$,$\alpha_1$,$\beta_2$,$\alpha_2$},
%     ytick=\empty
% ]
% \addplot[blue, thick, mark=none]
% table [x=s, y=PU] {figs/fig_parabola/parabola_case_b.txt};

% \addplot[red, thick, mark=none]
% table [x=s, y=PA] {figs/fig_parabola/parabola_case_b.txt};
% \end{axis}
% \end{tikzpicture}
% \caption{Attacker with medium variance.}
% \end{subfigure}
% \hfill
% \begin{subfigure}[t]{0.31\textwidth}
% \centering
% \begin{tikzpicture}
% \begin{axis}[
%     scale only axis,
%     width=0.8\linewidth,
%     height=0.45\linewidth,
%     xlabel={$s$},
%     ylabel={$P(s)$},
%     xmin=0.5, xmax=4.5,
%     ymin=-0.1, ymax=1.65,
%     xtick={0.6,3,3.5,4},
%     xticklabels={$\beta_1$,$\alpha_1$,$\beta_2$,$\alpha_2$},
%     ytick=\empty
% ]
% \addplot[blue, thick, mark=none]
% table [x=s, y=PU] {figs/fig_parabola/parabola_case_c.txt};

% \addplot[red, thick, mark=none]
% table [x=s, y=PA] {figs/fig_parabola/parabola_case_c.txt};
% \end{axis}
% \end{tikzpicture}
% \caption{Attacker with high variance.}
% \end{subfigure}

% \medskip

% \pgfplotslegendfromname{distlegend}

% \caption{Quadratic user and attacker distributions.}
% \label{fig:different_attackers_pdf_1_quadratic_cred}
% \end{figure*}

\begin{figure*}[t]
\centering

\begin{subfigure}[t]{0.31\textwidth}
\centering
\begin{tikzpicture}
\begin{axis}[
    scale only axis,
    width=0.8\linewidth,
    height=0.45\linewidth,
    xlabel={$t$},
    ylabel={$P_{\mathrm{success}}(t)$},
    xmin=1.5, xmax=4.5,
    ymin=0, ymax=1.05,
    xticklabels={$\beta_1$,$\alpha_1$,$\beta_2$,$\alpha_2$},
    xtick={2,3,3.5,4},
    ytick=\empty,
    legend to name=successlegend,
    legend columns=3,
    legend style={draw=black, fill=white}
]
\addplot[black, thick, mark=none]
table [x=s, y=Ps] {figs/fig_parabola/parabola_case_a.txt};
\addlegendentry{Success Curve}

\addplot[only marks, mark=*, mark size=1.3pt, blue]
table [x=T_opt, y=Ps_opt] {figs/fig_parabola/parabola_case_a_points.txt};
\addlegendentry{Optimal Point}

\addplot[only marks, mark=*, mark size=1.3pt, red]
table [x=T_eer, y=Ps_eer] {figs/fig_parabola/parabola_case_a_points.txt};
\addlegendentry{EER Point}
\end{axis}
\end{tikzpicture}
\caption{Attacker with low variance.}
\end{subfigure}
\hfill
\begin{subfigure}[t]{0.31\textwidth}
\centering
\begin{tikzpicture}
\begin{axis}[
    scale only axis,
    width=0.8\linewidth,
    height=0.45\linewidth,
    xlabel={$t$},
    ylabel={$P_{\mathrm{success}}(t)$},
    xticklabels={$\beta_1$,$\alpha_1$,$\beta_2$,$\alpha_2$},
    xtick={1.2,3,3.5,4},
    xmin=1, xmax=4.5,
    ymin=0, ymax=1.05,
    ytick=\empty
]
\addplot[black, thick, mark=none]
table [x=s, y=Ps] {figs/fig_parabola/parabola_case_b.txt};

\addplot[only marks, mark=*, mark size=1.3pt, blue]
table [x=T_opt, y=Ps_opt] {figs/fig_parabola/parabola_case_b_points.txt};

\addplot[only marks, mark=*, mark size=1.3pt, red]
table [x=T_eer, y=Ps_eer] {figs/fig_parabola/parabola_case_b_points.txt};
\end{axis}
\end{tikzpicture}
\caption{Attacker with medium variance.}
\end{subfigure}
\hfill
\begin{subfigure}[t]{0.31\textwidth}
\centering
\begin{tikzpicture}
\begin{axis}[
    scale only axis,
    width=0.8\linewidth,
    height=0.45\linewidth,
    xlabel={$t$},
    ylabel={$P_{\mathrm{success}}(t)$},
    xticklabels={$\beta_1$,$\alpha_1$,$\beta_2$,$\alpha_2$},
    xtick={0.6,3,3.5,4},
    xmin=0.5, xmax=4.5,
    ymin=0, ymax=1.05,
    ytick=\empty
]
\addplot[black, thick, mark=none]
table [x=s, y=Ps] {figs/fig_parabola/parabola_case_c.txt};

\addplot[only marks, mark=*, mark size=1.3pt, blue]
table [x=T_opt, y=Ps_opt] {figs/fig_parabola/parabola_case_c_points.txt};

\addplot[only marks, mark=*, mark size=1.3pt, red]
table [x=T_eer, y=Ps_eer] {figs/fig_parabola/parabola_case_c_points.txt};
\end{axis}
\end{tikzpicture}
\caption{Attacker with high variance.}
\end{subfigure}

\medskip

\pgfplotslegendfromname{successlegend}

\caption{Success probability for three quadratic attacker distributions.}
\label{fig:success_per_attacker_1_quadratic_cred}
\end{figure*}

% \begin{figure*}[t]
% \centering

% \begin{subfigure}[t]{0.31\textwidth}
% \centering
% \begin{tikzpicture}
% \begin{axis}[
%     scale only axis,
%     width=0.8\linewidth,
%     height=0.45\linewidth,
%     xlabel={$\FRR(t)$},
%     ylabel={$\FAR(t)$},
%     xmin=0, xmax=1,
%     ymin=0, ymax=1,
%     grid=none,
%     ytick =\empty,
%     xtick =\empty,
%     legend to name=farfrrlegend,
%     legend columns=3,
%     legend style={draw=black, fill=white}
% ]
% \addplot[thick, myPurple, mark=none]
% table [x=FRR, y=FAR]
% {figs/fig_parabola/parabola_case_a.txt};
% \addlegendentry{Operating Curve}

% \addplot[only marks, mark=*, mark size=1.3pt, blue]
% table [x=FRR_opt, y=FAR_opt]
% {figs/fig_parabola/parabola_case_a_points.txt};
% \addlegendentry{Optimal Point}

% \addplot[only marks, mark=*, mark size=1.3pt, red]
% table [x=FRR_eer, y=FAR_eer]
% {figs/fig_parabola/parabola_case_a_points.txt};
% \addlegendentry{\EER Point}
% \end{axis}
% \end{tikzpicture}
% \caption{Attacker with low variance.}
% \end{subfigure}
% \hfill
% \begin{subfigure}[t]{0.31\textwidth}
% \centering
% \begin{tikzpicture}
% \begin{axis}[
%     scale only axis,
%     width=0.8\linewidth,
%     height=0.45\linewidth,
%     xlabel={$\FRR(t)$},
%     ylabel={$\FAR(t)$},
%     xmin=0, xmax=1,
%     ymin=0, ymax=1,
%     grid=none,
%     ytick =\empty,
%     xtick =\empty
% ]
% \addplot[thick, myPurple, mark=none]
% table [x=FRR, y=FAR]
% {figs/fig_parabola/parabola_case_b.txt};

% \addplot[only marks, mark=*, mark size=1.3pt, blue]
% table [x=FRR_opt, y=FAR_opt]
% {figs/fig_parabola/parabola_case_b_points.txt};

% \addplot[only marks, mark=*, mark size=1.3pt, red]
% table [x=FRR_eer, y=FAR_eer]
% {figs/fig_parabola/parabola_case_b_points.txt};
% \end{axis}
% \end{tikzpicture}
% \caption{Attacker with medium variance.}
% \end{subfigure}
% \hfill
% \begin{subfigure}[t]{0.31\textwidth}
% \centering
% \begin{tikzpicture}
% \begin{axis}[
%     scale only axis,
%     width=0.8\linewidth,
%     height=0.45\linewidth,
%     xlabel={$\FRR(t)$},
%     ylabel={$\FAR(t)$},
%     xmin=0, xmax=1,
%     ymin=0, ymax=1,
%     grid=none,
%     ytick =\empty,
%     xtick =\empty
% ]
% \addplot[thick, myPurple, mark=none]
% table [x=FRR, y=FAR]
% {figs/fig_parabola/parabola_case_c.txt};

% \addplot[only marks, mark=*, mark size=1.3pt, blue]
% table [x=FRR_opt, y=FAR_opt]
% {figs/fig_parabola/parabola_case_c_points.txt};

% \addplot[only marks, mark=*, mark size=1.3pt, red]
% table [x=FRR_eer, y=FAR_eer]
% {figs/fig_parabola/parabola_case_c_points.txt};
% \end{axis}
% \end{tikzpicture}
% \caption{Attacker with high variance.}
% \end{subfigure}

% \medskip

% \pgfplotslegendfromname{farfrrlegend}

% \caption{\FAR vs. \FRR operating curves for different attacker distributions.}
% \label{fig:FAR_FRR_1_quadratic_cred}
% \end{figure*}

To demonstrate that this optimal threshold differs from the Equal Error Rate (\EER) threshold, we computed the~\EER threshold as follows:
\begin{equation*}
t_{\EER} = \frac{\alpha_1 \beta_1 - \alpha_2 \beta_2}{\alpha_1 - \alpha_2 + \beta_1 - \beta_2}
\end{equation*}
We can observe that the two thresholds differ.
To evaluate the difference between the~$t_{\text{opt}}$ and the~\EER, we consider different attacker distributions, each spanning a wider interval~$(\beta_1,\beta_2)$, which corresponds to increasing the variance.
Comparing each case with its corresponding success probability curve, we observe that the optimal threshold~$t_{\text{opt}}$ becomes more distinct from the~\EER threshold~$t_{\EER}$ as the variance of the attacker distribution increases and the more asymmetric the user and attacker distributions become. 

Moreover, increasing the variance of an attacker distribution shifts the optimal threshold~$t_{\text{opt}}$ towards~$\alpha_1$, reducing the~\FRR to zero. In contrast, the~\EER threshold remains fixed balancing between the~\FAR and~\FRR. This shift indicates that when increasing the variance, the optimal threshold prioritizes minimizing the false rejection rate, rather than balancing both error rates, as shown in Figure~\ref{fig:success_per_attacker_1_quadratic_cred}.

To quantify the effect of the variance on the deviation between the optimal threshold~$t_{\text{opt}}$ and~\EER threshold~$t_{\EER}$, we numerically computed this gap as a function of the attacker's variance. The simulation indicates that higher attacker variance enlarges the gap between the two success probabilities, as shown in Figure~\ref{fig:gap_vs_variance_parabola_success}. Operating at~$t_{\text{opt}}$ rather than the~\EER threshold is therefore most valuable when the attacker distribution spans a wide interval.

\begin{figure}[t]
\centering
\begin{tikzpicture}
\begin{axis}[
    scale only axis,
    width=0.82\columnwidth,
    height=0.48\columnwidth,
    xlabel={Attacker Variance $\mathrm{Var}(s)$},
    ylabel={$P_{\mathrm{success}}$},
    xtick=\empty,
    ytick=\empty,
    grid=none,
    legend style={
        at={(0.98,0.02)},
        anchor=south east,
        draw=black,
        fill=white
    }
]
\addplot[thick, blue, mark=none]
table [x=variance, y=Pmax] {figs/fig_parabola/gap_vs_variance_parabola.txt};
\addlegendentry{$P_{\max}$}

\addplot[thick, red, dashed, mark=none]
table [x=variance, y=PEER] {figs/fig_parabola/gap_vs_variance_parabola.txt};
\addlegendentry{$P_{\EER}$}
\end{axis}
\end{tikzpicture}

\caption{Comparison between the optimal threshold and the \EER threshold as a function of the attacker's variance.}
\label{fig:gap_vs_variance_parabola_success}
\end{figure}

\subsection{Empirical Evaluation}\label{sec:eval_single_cred}
\noindent
We first evaluate a mechanism that uses a biometric credential only under the baseline attack, where attackers authenticate using their own fingerprints. Calculating the mechanism’s success probability requires the biometric credential’s fault probabilities, which depend on the~\FAR and~\FRR. We therefore follow the evaluation procedure used by Jain et al.~\cite{Anil1997}, Cappelli et al.~\cite{Cappelli2006} and Slobodan et al.~\cite{Slob2005}. We produce histograms genuine and impostor comparisons, as shown in Figure~\ref{fig:sourceafis_histogram}and calculate the~\FAR and~\FRR as functions of the threshold~$t$, as shown in Figure~\ref{fig:approximations}.

Using these error rates, we calculate the biometric credential fault probabilities~(Equation~\ref{eqn:safe_prob}) and the mechanism's success probability~(Equation~\ref{eqn:safesuccess}). We find that the optimal threshold differs from the~\EER point and achieves a higher success probability, as shown in Figure~\ref{fig:bio_only_sourceafis_reg_scoring}. However, the success probability of the mechanism, when operated at the optimal threshold, is not significantly greater than operating at the~\EER threshold.

We next consider a more challenging scenario in which attackers use forged fingerprints to impersonate legitimate users. Unlike the baseline attack, spoofing samples can produce matching scores closer to those of genuine fingerprints, increasing the overlap between the distributions. In this evaluation, the~\EER increased significantly to 24.8\%, indicating that the biometric credential is less effective at distinguishing users from attackers. We therefore evaluate the mechanism's success probability under this attack, as shown in Figure~\ref{fig:livdet_sourceafis_success}. Nevertheless, despite the increased overlap between the user and the attacker distributions, the optimal and \EER thresholds differ, but the improvement in success probability remains small. Under the baseline attack, the success probability increases from 92.42\% to 92.54\%, while under spoofing attacks, it increases from 50.64\% to 50.88\%. Thus, the EER heuristic performs well for a single biometric credential mechanism in our evaluation.


\begin{figure*}[t]
\centering

\begin{subfigure}[t]{0.48\textwidth}
\centering
\begin{tikzpicture}
\begin{axis}[
    ybar,
    legend to name=histogramlegend,
    bar width=3pt,
    width=0.86\linewidth,
    height=0.40\linewidth,
    xlabel={Normalized similarity score (\%)},
    ylabel={Impostor count},
    ylabel style={red},
    grid=major,
    xmin=0, xmax=100,
    ymin=0, ymax=9000,
    enlarge x limits=false,
    enlarge y limits=false,
    tick style={draw=none},
    clip=true
]
\addplot+[ybar, fill=red, draw=black]
table [x=bin, y=impostor] {figs/fig_different_users_sourceafsi/sourceafis_histogram_data.txt};
\end{axis}

\begin{axis}[
    ybar,
    bar width=3pt,
    legend to name=histogramlegend,
    legend columns=2,
    width=0.86\linewidth,
    height=0.40\linewidth,
    xlabel={Normalized similarity score (\%)},
    yticklabel pos=right,
    ylabel={Genuine count},
    ylabel style={blue},
    grid=none,
    xmin=0, xmax=100,
    ymin=0, ymax=80,
    enlarge x limits=false,
    enlarge y limits=false,
    tick style={draw=none},
    clip=true
]
\addplot+[ybar, fill=blue, fill opacity=0.5]
table [x=bin, y=genuine] {figs/fig_different_users_sourceafsi/sourceafis_histogram_data.txt};
\addlegendentry{Genuine}

\addlegendimage{fill=red, mark=none}
\addlegendentry{Impostor}
\end{axis}
\end{tikzpicture}
\medskip
\pgfplotslegendfromname{histogramlegend}
\caption{Score histogram for genuine and impostor comparisons.}
\label{fig:sourceafis_histogram}
\end{subfigure}
\hfill
\begin{subfigure}[t]{0.48\textwidth}
\centering
\begin{tikzpicture}
\begin{axis}[
    scale only axis,
    width=0.82\linewidth,
    height=0.40\linewidth,
    xlabel={Threshold (\%)},
    ylabel={Success Probability},
    xmin=0, xmax=100,
    ymin=0, ymax=1,
    grid=major,
    legend style={
        at={(0.5,-0.30)},
        anchor=north,
        legend columns=3,
        draw=black,
        fill=white,
        font=\scriptsize
    },
    filter discard warning=false
]

\addplot[thick, blue, mark=none]
table [x=T, y=P_success] {figs/fig_different_users_sourceafsi/sourceafis_p_success_data.txt};
\addlegendentry{$P_{\mathrm{success}}(t)$}

\addplot[only marks, mark=square*, mark size=1.4pt, myGreen]
table [x=T_eer, y=P_eer] {figs/fig_different_users_sourceafsi/sourceafis_p_success_points.txt};
\addlegendentry{$P_{\mathrm{success}}(t_{\EER})$}

\addplot[only marks, mark=*, mark size=1.6pt, myRed]
table [x=T_opt, y=P_opt] {figs/fig_different_users_sourceafsi/sourceafis_p_success_points.txt};
\addlegendentry{Optimal point}

\end{axis}
\end{tikzpicture}
\caption{Single biometric credential mechanism's success.}
\label{fig:bio_only_sourceafis_reg_scoring}
\end{subfigure}

\caption{Baseline attack evaluation of a mechanism that uses a biometric credential only.}
\label{fig:sourceafis_success}
\end{figure*}

% ---------------------------------------------------------------------------
% Spoofing attack: score histogram and single-credential success.
% ---------------------------------------------------------------------------
\begin{figure*}[t]
\centering

\begin{subfigure}[t]{0.48\textwidth}
\centering
\begin{tikzpicture}
\begin{axis}[
    ybar,
    legend to name=livdethistogramlegend,
    bar width=2pt,
    width=0.86\linewidth,
    height=0.40\linewidth,
    xlabel={Normalized similarity score (\%)},
    ylabel={Impostor count},
    ylabel style={red},
    grid=major,
    xmin=0, xmax=100,
    ymin=0, ymax=600,
    enlarge x limits=false,
    tick style={draw=none},
    clip=true
]
\addplot+[ybar, fill=red, draw=black]
table[x=score, y=impostor]
{figs/fig_spoofed_users/livdet_sourceafis_histogram_best_of_probes_data.txt};
\end{axis}

\begin{axis}[
    ybar,
    legend to name=livdethistogramlegend,
    legend columns=2,
    bar width=2pt,
    width=0.86\linewidth,
    height=0.40\linewidth,
    xlabel={Normalized similarity score (\%)},
    ylabel={Genuine count},
    ylabel style={blue},
    yticklabel pos=right,
    grid=none,
    xmin=0, xmax=100,
    ymin=0, ymax=80,
    enlarge x limits=false,
    tick style={draw=none},
    clip=true
]
\addplot+[ybar, fill=blue, draw=blue, fill opacity=0.5]
table[x=score, y=genuine]
{figs/fig_spoofed_users/livdet_sourceafis_histogram_best_of_probes_data.txt};
\addlegendentry{Genuine}

\addlegendimage{fill=red, draw=black, mark=none}
\addlegendentry{Impostor}
\end{axis}
\end{tikzpicture}
\medskip
\pgfplotslegendfromname{livdethistogramlegend}
\caption{Score histogram for genuine and impostor comparisons.}
\label{fig:livdet_sourceafis_histogram}
\end{subfigure}
\hfill
\begin{subfigure}[t]{0.48\textwidth}
\centering
\begin{tikzpicture}
\begin{axis}[
    scale only axis,
    width=0.82\linewidth,
    height=0.40\linewidth,
    xlabel={Threshold (\%)},
    ylabel={Success Probability},
    xmin=0, xmax=100,
    ymin=0, ymax=1,
    grid=major,
    legend style={
        at={(0.5,-0.30)},
        anchor=north,
        legend columns=3,
        draw=black,
        fill=white,
        font=\scriptsize
    },
    each nth point=2,
    filter discard warning=false
]

\addplot[thick, blue, mark=none]
table [x=T, y=P_success] {figs/fig_spoofed_users/livdet_sourceafis_p_success_best_of_probes_data.txt};
\addlegendentry{$P_{\mathrm{success}}(t)$}

\addplot[only marks, mark=square*, mark size=1.4pt, myGreen]
table [x=T_eer, y=P_eer] {figs/fig_spoofed_users/livdet_sourceafis_p_success_best_of_probes_points.txt};
\addlegendentry{$P_{\mathrm{success}}(t_{\EER})$}

\addplot[only marks, mark=*, mark size=1.6pt, myRed]
table [x=T_opt, y=P_opt] {figs/fig_spoofed_users/livdet_sourceafis_p_success_best_of_probes_points.txt};
\addlegendentry{Optimal point}

\end{axis}
\end{tikzpicture}
\caption{Single biometric credential mechanism's success.}
\label{fig:livdet_sourceafis_success}
\end{subfigure}

\caption{Spoofing attack evaluation of a mechanism that uses a biometric credential only.}
\label{fig:livdet_single_cred_panel}
\end{figure*}


\section{Two Credentials}\label{sec:two_cred}
\noindent
We explore mechanisms that combine a biometric credential and an atomic credential. We symbolically analyze such mechanisms~$(\S\ref{sec:atm_biocred})$ and evaluate them on public datasets~$(\S\ref{sec:two_credentials_empirical_eval})$. We then characterize when such mechanisms achieve greater success than a mechanism using the atomic credential alone. We show they never do when the atomic credential can only be safe or stolen~$(\S\ref{sec:no_success_improvement})$. Finally, we present a necessary and sufficient condition for improvement~$(\S\ref{sec:necessary_sufficient_success_condition})$ and extend it to multiple atomic credentials.

\subsection{One Atomic and One Biometric Credential}\label{sec:atm_biocred}
\noindent
For a mechanism~$M$ combining two independently and identically distributed credentials, we defined mechanism's success under two scenarios:
a. \emph{AND}: The user can access both credentials, while the attacker cannot access at least one of the credentials. The mechanism succeeds when both credentials are in the safe state or when one is in the safe state and the other is in the leak state. Conversely, if either credential is in the lost or stolen state, the user cannot satisfy the mechanism, and the mechanism fails.
\begin{equation}
    \label{eqn:andsuccess}
P_{\text{Success}}(M^{\text{AND}}) = \mathrm{P^1_{safe}} \cdot \mathrm{P^2_{safe}} + \mathrm{P^1_{safe}} \cdot \mathrm{P^2_{leak}} + \mathrm{P^2_{safe}} \cdot \mathrm{P^1_{leak}}
\end{equation}
\noindent
b. \emph{OR}: The user can access at least one credential, while the attacker cannot access either of them. The mechanism succeeds when both credentials are in the safe state or when one is in the safe state and the other is in the lost state. If either credential is leaked or stolen, an attacker can satisfy the mechanism and thus the mechanism fails.
\begin{equation}
    \label{eqn:orsuccess}
P_{\text{Success}}(M^{\text{OR}}) = \mathrm{P^1_{safe}} \cdot \mathrm{P^2_{safe}} + \mathrm{P^1_{safe}} \cdot \mathrm{P^2_{loss}} + \mathrm{P^2_{safe}} \cdot \mathrm{P^1_{loss}}
\end{equation}

We start by analyzing the case of an \emph{AND} mechanism that combines an atomic credential and a biometric credential. 
We define the atomic credential with a constant value of each of its state probabilities as~$P^1_{safe}$, $P^1_{loss}$, $P^1_{leak}$, $P^1_{theft}$ corresponding respectively to each one of the states:~${\text{safe}}$,~${\text{loss}}$,~${\text{leak}}$ and~${\text{theft}}$. 
We define the biometric credential by its PDFs for the user and the attacker,~$f_U$ and~$f_A$ respectively. We consider uniform distribution as a first-order approximation of a biometric credential's distribution. Denote by~$u_1, u_2$ the user's PDF bounds and~$a_1, a_2$ the attacker's PDF bounds, as shown in Figure~\ref{fig:uniform_three_cases}. The state probabilities of the biometric credential are defined by the error metrics~\FAR and~\FRR as done in Section~\ref{sec:credentials}. Moreover, we consider the overlapping case between the user and attacker distributions, as this represents the common case among real-world scenarios. Thus, we assume that~$a_1 < u_1 < a_2 < u_2$.
We set the state probabilities of an atomic and a biometric credential into Equation~\ref{eqn:andsuccess}. Thus, the probability of success is a function of the threshold $t$ and divided into 5 intervals. Denote:
\[L(t) = \dfrac{u_2 - t}{u_2 - u_1}
\left[
P^{1}_{\mathrm{safe}} + P^{1}_{\mathrm{leak}}(1
- \dfrac{a_2 - t}{a_2 - a_1})
\right].\]
\[
P_{\mathrm{success}}(t)=
\begin{cases}
P^{1}_{\mathrm{safe}}, 
& t \le a_1 \\

P^{1}_{\mathrm{safe}} + P^{1}_{\mathrm{leak}}(1-\dfrac{a_2 - t}{a_2 - a_1}), 
& a_1 < t \le u_1 \\

L(t),
& u_1 < t \le a_2 \\

\dfrac{u_2 - t}{u_2 - u_1}\cdot(P^{1}_{\mathrm{safe}} + P^{1}_{\mathrm{leak}}),
& a_2 < t \le u_2 \\

0, 
& t > u_2 .
\end{cases}
\]

The success probability is constant in the first interval, strictly increasing in the second, quadratic in the third, strictly decreasing in the fourth, and zero in the fifth.
The global optimum therefore lies in the third interval, $u_1 < t \le a_2$, where the user and attacker distributions overlap and the success function attains an interior maximum:
\begin{equation*}
t_{\text{opt}}^{\mathrm{AND}}
=
\frac{
u_2 + a_2
- (a_2 - a_1)\,\dfrac{P^{1}_{\text{safe}} + P^{1}_{\text{leak}}}{P^{1}_{\text{leak}}}
}{2},
\end{equation*}
clipped to the interval boundaries.
Evaluating at the lower boundary gives $P_{\text{success}}(u_1) > P^{1}_{\text{safe}}$, so integrating the biometric credential strictly improves the mechanism over using the atomic credential alone.

The optimum's position reveals which credential the mechanism relies on.
As the variance of the attacker distribution increases,~$t_{\text{opt}}^{\mathrm{AND}}$ shifts toward the lower boundary~$u_1$.
An atomic credential with high~$P^{1}_{\text{safe}}$ and low~$P^{1}_{\text{leak}}$ shifts it toward $u_1$ as well, and the mechanism's success then depends mainly on the atomic credential.
Conversely, high~$P^{1}_{\text{leak}}$ and low~$P^{1}_{\text{safe}}$ shift it toward the upper boundary~$a_2$, and the mechanism relies more heavily on the biometric credential.

We obtain a similar analysis for the \emph{OR} mechanism, with~$P^{1}_{\mathrm{loss}}$ in place of~$P^{1}_{\mathrm{leak}}$.
Its global optimum lies in the same third interval:
\begin{equation*}
t_{\text{opt}}^{\mathrm{OR}} = \frac{a_1 P^{1}_{\mathrm{loss}} + u_2 (P^{1}_{\mathrm{safe}} + P^{1}_{\mathrm{loss}}) - P^{1}_{\mathrm{safe}} u_1}{2 P^{1}_{\mathrm{loss}}}.
\end{equation*}
For both mechanisms the optimal threshold differs from the~\EER threshold,~$t_{\EER} = (a_1 u_1 - a_2 u_2)/(a_1 - a_2 + u_1 - u_2)$.

We observe that the optimal threshold of the biometric credential in both mechanisms is not, in general, the same as its optimal threshold when used without the atomic credential.
An integrator must therefore set the threshold with respect to the mechanism it is deployed in, and not treat the biometric credential in isolation.
The full interval analysis and the resulting success probabilities are in the full version of this report~\cite{biometric_full_version}.

\subsection{Empirical Evaluation}\label{sec:two_credentials_empirical_eval}
\noindent
We use the same biometric credential evaluated in~$(\S\ref{sec:eval_single_cred})$, and first consider the baseline attack scenario, where attackers use their own fingerprints to authenticate. The mechanisms combine this biometric credential with an atomic credential whose fault probabilities are arbitrarily set to~$P^{1}_\text{safe} = 0.75$, $P^{1}_\text{leak} = 0.1$, $P^{1}_\text{loss} = 0.1$, and $P^{1}_\text{stolen} = 0.05$. We then calculate the success probabilities of the \emph{AND mechanism}~(Equation~\ref{eqn:andsuccess}), and the \emph{OR mechanism}~(Equation~\ref{eqn:orsuccess}). We illustrate these probabilities as functions of the threshold~$t$ in Figure~\ref{fig:bio_atm_sourceafis_reg_scoring}. We show that the optimal threshold for each mechanism differs from the~\EER threshold. Moreover, both mechanisms operating at their optimal thresholds achieve success probabilities that are~8-9\% higher than that of a mechanism using only the atomic credential, demonstrating the benefit of combining the two credentials.

However, in this attack scenario, optimizing the thresholds provides only a small improvement over the \EER threshold for both mechanisms, since the biometric credential already has a very low \EER, indicating limited overlap between the genuine and impostor distributions.

We next evaluate the two mechanisms under spoofing attacks, where attackers use forged fingerprints to impersonate legitimate users. As shown in the previous subsection, spoofing increases the overlap between the genuine and impostor score distributions, making authentication more challenging. We calculate the success probabilities of both mechanisms under this attack, as shown in Figure~\ref{fig:livdet_and_or_success}. We find that the optimal threshold reduces the failure probability between~20\%-72\% compared to the~\EER threshold for the two mechanisms. Moreover, a mechanism combining both credentials at the~\EER threshold may achieve a lower success probability than a mechanism using either credential alone. These results demonstrate that using a non-optimal biometric threshold can reduce the success probability when combining credentials in an authentication mechanism. The results are summarized in Table~\ref{tab:psafe_sweep_improvement}.

Since threshold selection can either reduce or improve the mechanism's success, the threshold of the biometric credential must be selected in the context of the mechanism.
This motivates us to characterize when combining a biometric credential with an atomic credential improves the mechanism at all.

% ===========================================================================
% FIGURES AND TABLE FOR sec:two_credentials_empirical_eval
% ===========================================================================

% ---------------------------------------------------------------------------
% Baseline attack: AND and OR mechanisms.
% ---------------------------------------------------------------------------
\begin{figure}[t]
\centering
\begin{tikzpicture}
\begin{axis}[
    scale only axis,
    width=0.82\columnwidth,
    height=0.36\columnwidth,
    xlabel={Threshold (\%)},
    ylabel={Success Probability},
    xmin=0, xmax=100,
    ymin=0, ymax=1,
    grid=major,
    legend style={
        at={(0.5,-0.3)},
        anchor=north,
        legend columns=3,
        draw=black,
        fill=white,
        font=\scriptsize
    },
    filter discard warning=false
]

\addplot[thick, blue, mark=none]
table [x=T, y=P_and] {figs/fig_different_users_sourceafsi/sourceafis_success_and_or_data.txt};
\addlegendentry{$P_{\mathrm{success}}^{\mathrm{AND}}(t)$}

\addplot[thick, myOrange, mark=none]
table [x=T, y=P_or] {figs/fig_different_users_sourceafsi/sourceafis_success_and_or_data.txt};
\addlegendentry{$P_{\mathrm{success}}^{\mathrm{OR}}(t)$}

\addplot[only marks, mark=*, mark size=1.5pt, blue]
table [x=T_and_opt, y=P_and_opt] {figs/fig_different_users_sourceafsi/sourceafis_success_and_or_points.txt};
\addlegendentry{AND optimum}

\addplot[only marks, mark=x, mark size=1.9pt, myOrange]
table [x=T_or_opt, y=P_or_opt] {figs/fig_different_users_sourceafsi/sourceafis_success_and_or_points.txt};
\addlegendentry{OR optimum}

\addplot[only marks, mark=square*, mark size=1.4pt, myGreen]
table [x=T_eer, y=P_and_eer] {figs/fig_different_users_sourceafsi/sourceafis_success_and_or_points.txt};
\addlegendentry{AND at \EER}

\addplot[only marks, mark=diamond*, mark size=1.5pt, myRed]
table [x=T_eer, y=P_or_eer] {figs/fig_different_users_sourceafsi/sourceafis_success_and_or_points.txt};
\addlegendentry{OR at \EER}

\addplot[thick, black, dashed, domain=0:100] {0.75};
\addlegendentry{Atomic credential only}

\end{axis}
\end{tikzpicture}
\caption{Baseline attack. A biometric credential integrated with an atomic credential into AND and OR mechanisms. The AND and OR mechanisms achieve success probabilities of $0.78$ and $0.82$ at the \EER point, and $0.81$ and $0.84$ at their respective optimal operating points.}
\label{fig:bio_atm_sourceafis_reg_scoring}
\end{figure}

% ---------------------------------------------------------------------------
% Spoofing attack: AND and OR mechanisms.
% ---------------------------------------------------------------------------
\begin{figure}[t]
\centering
\begin{tikzpicture}
\begin{axis}[
    scale only axis,
    width=0.82\columnwidth,
    height=0.36\columnwidth,
    xlabel={Threshold (\%)},
    ylabel={Success Probability},
    xmin=0, xmax=100,
    ymin=0, ymax=1,
    grid=major,
    legend style={
        at={(0.5,-0.3)},
        anchor=north,
        legend columns=3,
        draw=black,
        fill=white,
        font=\scriptsize
    },
    each nth point=2,
    filter discard warning=false
]

\addplot[thick, blue, mark=none]
table [x=T, y=P_and] {figs/fig_spoofed_users/livdet_sourceafis_success_and_or_best_of_probes_data.txt};
\addlegendentry{$P_{\mathrm{success}}^{\mathrm{AND}}(t)$}

\addplot[thick, myOrange, mark=none]
table [x=T, y=P_or] {figs/fig_spoofed_users/livdet_sourceafis_success_and_or_best_of_probes_data.txt};
\addlegendentry{$P_{\mathrm{success}}^{\mathrm{OR}}(t)$}

\addplot[only marks, mark=*, mark size=1.5pt, blue]
table [x=T_and_opt, y=P_and_opt] {figs/fig_spoofed_users/livdet_sourceafis_success_and_or_best_of_probes_points.txt};
\addlegendentry{AND optimal point}

\addplot[only marks, mark=x, mark size=1.9pt, myOrange]
table [x=T_or_opt, y=P_or_opt] {figs/fig_spoofed_users/livdet_sourceafis_success_and_or_best_of_probes_points.txt};
\addlegendentry{OR optimal point}

\addplot[only marks, mark=square*, mark size=1.4pt, myGreen]
table [x=T_eer, y=P_and_eer] {figs/fig_spoofed_users/livdet_sourceafis_success_and_or_best_of_probes_points.txt};
\addlegendentry{AND at \EER}

\addplot[only marks, mark=diamond*, mark size=1.5pt, myRed]
table [x=T_eer, y=P_or_eer] {figs/fig_spoofed_users/livdet_sourceafis_success_and_or_best_of_probes_points.txt};
\addlegendentry{OR at \EER}

\addplot[thick, black, dashed, domain=0:100] {0.8};
\addlegendentry{Atomic credential only}

\end{axis}
\end{tikzpicture}
\caption{Spoofing attack. A biometric credential integrated with an atomic credential into AND and OR mechanisms. The AND mechanism achieves a maximal success probability of $0.834$, while the OR mechanism achieves $0.816$. At the \EER operating point, the success probabilities are $0.713$ and $0.649$ for the AND and OR mechanisms, respectively.}
\label{fig:livdet_and_or_success}
\end{figure}

% ---------------------------------------------------------------------------
% Spoofing attack: success probability sweep over atomic state probabilities.
% ---------------------------------------------------------------------------
\begin{table*}[t]
\centering
\caption{Comparison between the success probabilities obtained at the EER threshold and the maximal success probabilities. Bold EER values indicate the worse composition, whereas bold maximal values indicate the best composition. The bold safe column indicating the success probability of a mechanism uses an atomic credential only. The improvement column shows the percentage of improvement between the failure probability at the optimal mechanism success and the failure probability at the maximal mechanism success at EER threshold.}
\label{tab:psafe_sweep_improvement}
\begin{tabular}{cccccccccc}
\toprule
$\mathbf{P_{\mathrm{safe}}}$ & $P_{\mathrm{loss}}$ & $P_{\mathrm{leak}}$ & $P_{\mathrm{theft}}$ & AND at EER & OR at EER & Max AND & Max OR & Improvement \\
\midrule
\textbf{0.60} & 0.30 & 0.05 & 0.05 & \textbf{0.45} & 0.58 & 0.60 & \textbf{0.66} & 20.25\% \\
\textbf{0.60} & 0.05 & 0.30 & 0.05 & 0.57 & \textbf{0.45} & \textbf{0.68} & 0.60 & 24.94\% \\
\textbf{0.75} & 0.15 & 0.05 & 0.05 & \textbf{0.56} & 0.61 & 0.75 & \textbf{0.77} & 40.39\% \\
\textbf{0.82} & 0.05 & 0.08 & 0.05 & 0.62 & \textbf{0.61} & \textbf{0.824} & 0.82 & 53.20\% \\
\textbf{0.90} & 0.00 & 0.05 & 0.05 & 0.66 & \textbf{0.64} & 0.89 & \textbf{0.905} & 72.19\% \\
\bottomrule
\end{tabular}
\end{table*}

\subsection{No Improvement In Mechanism Success}\label{sec:no_success_improvement}
\noindent
We first consider an atomic credential that is always either safe or stolen, never lost or leaked. This case admits no benefit from producing a mechanism that uses an atomic credential in this setting and a biometric credential. The following proposition formalizes this result.
\begin{proposition}
    Let~$c_1$ be an atomic credential with~${P^1_\text{loss} = P^1_\text{leak} = 0}$, $c_2$ be a biometric credential with user distribution~$f_U(t)$ and attacker distribution~$f_A(t)$, and~$M$ be a mechanism using~$c_1$. Then, for a mechanism~$\bar{M}$ using~$c_1$,~$c_2$ in \emph{AND} or \emph{OR} compositions it holds that
    \begin{equation*}
        P_{\mathrm{success}}^{\bar{M}} \le P_{\mathrm{success}}^{M}.
    \end{equation*}
\end{proposition}
\noindent
Intuitively, a mechanism that uses the atomic credential alone succeeds whenever that credential is safe. 
Adding a biometric credential can only reject the user in some of those cases. 
It never helps when the atomic credential is stolen, since the attacker then satisfies the atomic credential regardless of the biometric one. 
The proof is in the full version of this report~\cite{biometric_full_version}.

\subsection{Necessary and Sufficient Condition}\label{sec:necessary_sufficient_success_condition}
\noindent 
We turn to the general case of an atomic credential that can also be lost or leaked. Then, we identify a biometric credential that is \emph{asymptotically separable}. We show that this is a necessary and sufficient condition for a mechanism using an atomic credential and a biometric credential to achieve strictly greater success probability than a mechanism using only the atomic credential.
Intuitively, asymptotic separability means that there exists a threshold~$t$ such that the error rate is negligible with respect to the correctness rate, making the user and attacker distributions effectively separated.
\begin{definition}
A user distribution function~$f_U(s)$ and an attacker distribution function~$f_A(s)$ are \emph{asymptotically separable} if
\begin{equation*}
    \inf_{t}
    \min\!\left\{
    \frac{\FAR(t)}{1-\FRR(t)},\,
    \frac{\FRR(t)}{1-\FAR(t)}
    \right\}
    = 0.
\end{equation*}
\end{definition}
We now prove that asymptotic separability of a biometric credential is a sufficient condition.
\begin{proposition}\label{thm:suf_nec_cond_for_improve}
    Let~$c_1$ be an atomic credential with~$P^1_\text{safe} > 0$ and at least one of~$P^1_\text{leak},P^1_\text{loss} > 0$, $c_2$ be a biometric credential whose user distribution~$f_U(t)$ and attacker distribution~$f_A(t)$ are asymptotically separable and~$M$ be a mechanism using~$c_1$. Then, there exists a mechanism~$\bar{M}$ which is an \emph{AND} or \emph{OR} composition using~$c_1$ and~$c_2$, such that
    \begin{equation*}
        P_{\mathrm{success}}^{\bar{M}} > P_{\mathrm{success}}^{M}
    \end{equation*}
\end{proposition}
\noindent
The \emph{AND} mechanism improves on~$M$ exactly when the ratio~$\FRR(t)/(1-\FAR(t))$ falls below a positive bound determined by the atomic credential’s fault probabilities, and the \emph{OR} mechanism when~$\FAR(t)/(1-\FRR(t))$ does.
Both bounds are strictly positive whenever the atomic credential has $P^1_\text{safe} > 0$ and at least one of $P^1_\text{leak}, P^1_\text{loss} > 0$.
Asymptotic separability states that one of these two ratios comes arbitrarily close to zero, so a threshold satisfying the corresponding bound exists. The vanishing ratio selects the composition.
The proof is in the full version of this report~\cite{biometric_full_version}.

We showed that \emph{asymptotic separability} is a sufficient condition. We now show it is a necessary condition for producing a new mechanism~$\bar{M}$ with greater success probability than that of~$M$.
\begin{proposition}\label{thm:nec_cond_for_improve}
    Let~$c_2$ be a biometric credential with user distribution~$f_U(t)$ and attacker distribution~$f_A(t)$. If for every atomic credential~$c_1$ with~$P^1_\text{safe} > 0$ and at least one of~$P^1_\text{leak},P^1_\text{loss} > 0$, there exists a threshold~$t$ and a mechanism~$\bar{M}$ formed as an \emph{AND} or \emph{OR} composition such that
    \begin{equation*}
    P^{\bar{M}}_{\text{success}}(t) > P^{M}_{\text{success}}.
    \end{equation*}
    Then the biometric credential~$c_2$ is asymptotically separable.
\end{proposition} 
\noindent
Suppose the biometric credential is not asymptotically separable, so both~$\FRR(t)/(1-\FAR(t))$ and~$\FAR(t)/(1-\FRR(t))$ stay above some $\varepsilon > 0$ for every threshold.
The hypothesis quantifies over \emph{every} atomic credential, so it must hold in particular for one whose leak and loss probabilities are small enough relative to its safe probability that both bounds of Proposition~\ref{thm:suf_nec_cond_for_improve} fall below~$\varepsilon$.
No threshold can then satisfy either bound, so no composition improves on~$M$ for that credential, contradicting the hypothesis.
The proof is in the full version of this report~\cite{biometric_full_version}.

We showed that~$\bar{M}$ achieves greater success probability than~$M$ when the biometric credential is \emph{asymptotically separable}, and the atomic credential has~$P^1_\text{safe} > 0$ and at least one of~$P^1_\text{leak},P^1_\text{loss} > 0$.
We also showed that asymptotic separability is necessary for~$\bar{M}$ to achieve greater success probability than every atomic credential mechanism.
We extend our proof to a mechanism that uses a biometric credential and~$k$ atomic credentials, and show that it achieves greater success probability than a mechanism using only the atomic credentials.

First, every authentication mechanism~$M$ can be represented as a sum of products, i.e., a disjunctive normal form~(DNF)~\cite{Mendelson1997}. Then, we define a \emph{consolidated credential} of~$k$ atomic credentials as if it were itself a single atomic credential.

\begin{definition}
Let~$M$ be an authentication mechanism over~$k$ atomic credentials. The \emph{consolidated credential} of~$M$ is the single abstract credential whose availability to the user and attacker is determined by the outputs of~$M$ on their respective availability vectors~$\sigma^U,\sigma^A\in\{0,1\}^k$. The fault probabilities of the consolidated credential are thus
\begin{align*}
P^M_{\text{safe}}
&=
\Pr\left[M(\sigma^U)=1,\ M(\sigma^A)=0\right],\\
P^M_{\text{loss}}
&=
\Pr\left[M(\sigma^U)=0,\ M(\sigma^A)=0\right],\\
P^M_{\text{leak}}
&=
\Pr\left[M(\sigma^U)=1,\ M(\sigma^A)=1\right],\\
P^M_{\text{stolen}}
&=
\Pr\left[M(\sigma^U)=0,\ M(\sigma^A)=1\right].
\end{align*}
\end{definition}
\begin{proposition}\label{thm:suf_nec_cond_for_improve_k_atm}
    Let~$\{c_1,c_2,\ldots,c_k\}$ be a set of~$k$ independent atomic credentials,~$c_{k+1}$ be a biometric credential whose user distribution~$f_U(t)$ and attacker distribution~$f_A(t)$ are asymptotically separable and~$M$ be a mechanism using~$\{c_1,c_2,\ldots,c_k\}$ whose consolidated credential has~$P^M_\text{safe} > 0$ and at least one of~$P^M_\text{leak},P^M_\text{loss} > 0$. Then, there exists a mechanism~$\bar{M}$ using~$\{c_1,c_2,\ldots,c_k,c_{k+1}\}$ such that
    \begin{equation*}
        P_{\mathrm{success}}^{\bar{M}} > P_{\mathrm{success}}^{M}
    \end{equation*}
\end{proposition}
\begin{proof}
    Let~$M$ be a mechanism using~$\{c_1,c_2,\ldots,c_k\}$. The mechanism~$M$ induces a reduction of the~$k$ credentials into a single consolidated credential. The consolidated credential of~$M$ is a single atomic credential with state probabilities~$P^M_\text{safe}>0$, and at least one of~$P^M_\text{leak}$, or~$P^M_\text{loss} > 0$. By Proposition~\ref{thm:suf_nec_cond_for_improve}, there exists a mechanism~$\bar{M}$ using the consolidated credential and the biometric credential that achieves greater success probability than that of the consolidated credential mechanism. Since the consolidated credential is equivalent to the original mechanism~$M$, it follows that there exists a mechanism~$\bar{M}$ using~$\{c_1,c_2,\ldots,c_k,c_{k+1}\}$ that achieves greater success probability than that of~$M$.
\end{proof}
\begin{proposition}
    Let~$c_{k+1}$ be a biometric credential with user distribution~$f_U(t)$ and attacker distribution~$f_A(t)$. Suppose that for every set of~$k$ independent atomic credentials~$\{c_1,\ldots,c_k\}$ and every mechanism~$M$ over these credentials whose consolidated credential has~$P^M_\text{safe} > 0$ and at least one of~$P^M_\text{leak},P^M_\text{loss} > 0$, there exists a mechanism~$\bar{M}$ using~$\{c_1,\ldots,c_k,c_{k+1}\}$ such that
    \begin{equation*}
    P^{\bar{M}}_{\text{success}}(t) > P^{M}_{\text{success}}.
    \end{equation*}
    Then the biometric credential~$c_{k+1}$ must be asymptotically separable.
\end{proposition} 
\begin{proof}
    The mechanism~$M$ induces a reduction of the~$k$ credentials into a single consolidated credential. The consolidated credential of~$M$ is a single atomic credential with state probabilities~$P^M_\text{safe}>0$, and at least one of~$P^M_\text{leak},P^M_\text{loss} > 0$. Then,~$\bar{M}$ is a mechanism using the consolidated credential of~$M$ and the biometric credential. By Proposition~\ref{thm:nec_cond_for_improve}, the biometric credential must be asymptotically separable.
\end{proof}

Since biometric credentials often exhibit a more common behavior~\cite{Slob2005,anil2004}, we introduce the stronger \emph{asymptotic dominance} condition to capture this property. Intuitively, asymptotic dominance means that one distribution function dominates the other in the tails of the distributions: for very low score values~$s$ the attacker distribution dominates the user distribution, and vice versa for very high values of~$s$. We prove that a mechanism~$\bar{M}$ that uses an \emph{asymptotically dominated} biometric credential and an atomic credential, achieves a greater success probability than that of a single atomic credential mechanism~$M$. We prove that this is a sufficient condition~(see Appendix~\ref{app:sufficient-condition}). 

\section{Multiple Attempts}
\label{sec:multiple_auth_attempts}
\noindent
In biometric real world authentication, a mechanism commonly allows the principal to retry following failure~\cite{Chen2023BrutePrint}, typically up to a bounded number of attempts before lockout~\cite{typical_aut_attempts}. Since retries change the mechanism and raise the success probability of users and attackers alike, the single attempt threshold is no longer optimal. We therefore formalize a mechanism that limits the number of attempts~$(\S\ref{sec:equal_k_attempts})$. Within that limit, the mechanism success probability alone is not enough, since the threshold optimization may require many attempts to succeed. Then, we analyze both the mechanism's success probability and the mean number of attempts required for a user's first successful authentication~$(\S\ref{sec:mul_attempts_success})$. As these two metrics conflict, we examine their tradeoff subject to different constraints~$(\S\ref{sec:optimization_mul_attempts})$.

\subsection{Predefined Numbers of Attempts}\label{sec:equal_k_attempts}
\noindent
A deployed mechanism fixes its upper attempts number bound before any authentication takes place.
We first consider a mechanism that allows at most~$k$ attempts, and show how changing the predefined attempts number~$k$ affects the mechanism success probability. The mechnanism may use a different threshold~$t_j$ at every attempt~$j$, and the user fails to authenticate if he is rejected in all~$k$ attempts. Thus, the probability that the user is rejected at an arbitrary single attempt~$j$ is~$\FRR(t_j)$ and the probability that the user is accepted at least once, is the probability it's not rejected in all~$k$ attempts:~$1-\prod_{j=1}^{k}\FRR(t_j).$
Mechanism success requires rejecting the attacker at every attempt. Since the attacker is rejected at each attempt~$j$ with probability~$1-\FAR(t_j)$, the probability that the attacker is rejected in all~$k$ attempts is:~$\prod_{j=1}^{k}\left(1-\FAR(t_j)\right).$ The authentication mechanism succeeds when the user is accepted at least once and the attacker is rejected in every attempt. Therefore, the overall success probability is

\begin{equation}\label{eqn:success_multiple_attempts}
P_{\mathrm{success}}^{(k)}(\mathbf{t})
=
\left(
1-
\prod_{j=1}^{k}\FRR(t_j)
\right)
\prod_{j=1}^{k}
\left(1-\FAR(t_j)\right).
\end{equation}
Both the~$\FRR(t_j)$ and $\FAR(t_j)$ error rates are, by definition, lower or equal than~1, for every attempt~$j$ and threshold~$t_j$. The products of~$\FRR(t_j)$ and the products of $1-\FAR(t_j)$ are both monotonically decreasing functions with~$k$. As a result, for a fixed thresholds selection, increasing the attempts number~$k$ increases the user probability to be accepted, while reducing the probability of the attacker to be rejected. Increasing~$k$ can therefore help the user gain access, but also increases the risk of accepting an attacker. The thresholds selection, when optimizing the success probability, must therefore balance these two effects.

\subsection{User Authentication Within~$k$ Attempts}\label{sec:mul_attempts_success}
\noindent
We are now interested not only in whether the user eventually authenticates, but also in how many attempts it takes for a user successful authentication. We proceed in three steps. First, we derive the probability that the user's first successful authentication occurs at a given attempt, which yields the probability of authenticating within~$k$ attempts. Then, we consider the conditional attempt number given that the user successfully authenticates. Finally, we take the mean of this conditional distribution, obtaining the average number of attempts a user requires.
Let~$N_U$ denote the attempt number at which the user authenticates for the first time.
If the mechanism rejects the user in all~$k$ attempts, we define~$N_U =\infty$.
The event~$\{N_U=j\}$ means that the user is rejected in attempts~$1,\ldots,j-1$ and is then accepted at attempt~$j$.
Under the assumption that every authentication attempt is independent, the probability that the user authenticates for the first time at attempt~$j$ is therefore

\begin{equation}\label{eq:fist_success}
\Pr[N_U=j]
=
\left(
\prod_{i=1}^{j-1}\FRR(t_i)
\right)
\left(1-\FRR(t_j)\right).
\end{equation}
This is the \emph{first success probability} that the user successfully authenticates specifically at attempt~$j$, having failed every preceding attempt.
Summing the first success probabilities over all permitted attempts gives the probability that the user successfully authenticates within~$k$ attempts: 

\begin{equation*}
\Pr[N_U\leq k]
=
\sum_{j=1}^{k}\Pr[N_U=j].
\end{equation*}
Substituting the first success probabilities yields
\begin{equation*}
\Pr[N_U\leq k]
=
\sum_{j=1}^{k}
\left(
\prod_{i=1}^{j-1}\FRR(t_i)
\right)
\left(1-\FRR(t_j)\right).
\end{equation*}
As described in the previous section, this probability equals
\begin{equation}\label{eq:auth_within_ku}
\Pr[N_U\leq k]
=
1-\prod_{i=1}^{k}\FRR(t_i).
\end{equation}

Having derived the probability that the user successfully authenticates within~$k$ attempts, we now consider the conditional distribution of the attempt number~$N_U$ given that the user successfully authenticates within~$k$ attempts.
By the definition of conditional probability,

\begin{equation*}
\Pr[N_U=j\mid N_U\leq k]
=
\frac{
\Pr[
\{N_U=j\}
\cap
\{N_U\leq k\}
]
}{
\Pr[N_U\leq k]
}.
\end{equation*}
For every~$j\in\{1,\ldots,k\}$, the event~$\{N_U=j\}$ implies that the user succeeds within~$k$ attempts. In set notation,

\begin{equation*}
\{N_U=j\}
\subseteq
\{N_U\leq k\}.
\end{equation*}
Consequently, their intersection is simply
\begin{equation*}
\{N_U=j\}
\cap
\{N_U\leq k\}
=
\{N_U=j\}.
\end{equation*}
Then it follows that
\begin{equation*}
\Pr[N_U=j\mid N_U\leq k]
=
\frac{
\Pr[N_U=j]
}{
\Pr[N_U\leq k]
}.
\end{equation*}
Using Equation~\ref{eq:fist_success} and Equation~\ref{eq:auth_within_ku}, we obtain
\begin{equation*}
\Pr[N_U=j\mid N_U\leq k]
=
\frac{
\left(
\prod_{i=1}^{j-1}\FRR(t_i)
\right)
\left(1-\FRR(t_j)\right)
}{
1-
\prod_{i=1}^{k}\FRR(t_i)
}.
\end{equation*}
The numerator is the probability that the first successful authentication occurs exactly at attempt~$j$. The denominator is the probability that the user succeeds at some attempt no later than~$k$. Their ratio therefore gives the fraction of successful users whose first successful authentication occurs at attempt~$j$.

The conditional mean number of attempts is the expected value of~$N_U$ among users who successfully authenticate within~$k$ attempts. By the definition of conditional expectation for a discrete random variable,
\begin{equation*}
\mathbb{E}
\left[
N_U
\mid
N_U\leq k
\right]
=
\sum_{j=1}^{k}
j
\Pr[N_U=j\mid N_U\leq k].
\end{equation*}
Substituting the conditional probabilities gives

\begin{equation*}
\mathbb{E}
\left[
N_U
\mid
N_U\leq k
\right]
=
\frac{
\displaystyle
\sum_{j=1}^{k}
j
\left(
\prod_{i=1}^{j-1}\FRR(t_i)
\right)
\left(1-\FRR(t_j)\right)
}{
\displaystyle
1-
\prod_{j=1}^{k}\FRR(t_j)
}.
\end{equation*}
Conditioning on~$N_U\leq k$ is necessary because~$N_U=\infty$ when the user is rejected in all permitted attempts. The conditional mean measures the average number of attempts required specifically among successful authentication sessions, helping system designers to determine their requirements and priorities when optimizing thresholds over attempts.

\subsection{Optimization Under Mean-Success Trade-Off}\label{sec:optimization_mul_attempts}
\noindent
System designers evaluate the security of a multiple attempt authentication mechanism by two metrics: (1)~the conditional mean number of a user authentication attempts, and (2)~the overall success probability of the mechanism. The first metric measures usability, while the second metric measures reliability. However, there is an inherent trade-off between these metrics. As the number of attempts grows, the overall success probability increases while usability degrades. 

We consider two optimization problems: one that minimizes the conditional mean number of user authentication attempts subject to a constraint on the success probability, and another that maximizes the success probability subject to a constraint on the conditional mean number of user authentication attempts.  

We start by minimizing the conditional mean number of user authentication attempts. Thus
\begin{equation*}
\begin{aligned}
\underset{\mathbf{t}}{\operatorname{minimize}}
\quad&
\mathbb{E}
\left[
N_U
\mid
N_U\leq k
\right]
\\
\text{subject to}
\quad&
P_{\mathrm{success}}(\mathbf{t})
\geq
P'.
\end{aligned}
\end{equation*}
Equivalently, the constraint is

\begin{equation*}
\left(
1-
\prod_{j=1}^{k}\FRR(t_j)
\right)
\prod_{j=1}^{k}
\left(1-\FAR(t_j)\right)
\geq
P'.
\end{equation*}
Without the constraint, this problem has a trivial solution: set every threshold to accept everything, and the user always succeeds on the first attempt. The mean is then~1, but so is the attacker’s success rate. The constraint~$P'$ rules this out. At the optimum, the constraint holds with equality. Any solution with a strictly higher success probability could trade that slack for a more permissive threshold and a lower mean. Thus, this optimization enforces a minimum security level while maximizing the usability within it. 

We next maximize the mechanism's success probability subject to an upper
bound on the conditional mean number of user authentication attempts.
Let~$\mu'$ denote this upper bound. For fixed attempt limits~$k$, we optimize the threshold vector~$\mathbf{t}=(t_1,\ldots,t_k)$. Thus
\begin{equation*}
\begin{aligned}
\underset{\mathbf{t}}{\operatorname{maximize}}
\quad&
P_{\mathrm{success}}(\mathbf{t})
\\
\text{subject to}
\quad&
\mathbb{E}\left[N_U\mid N_U\leq k\right]
\leq \mu'.
\end{aligned}
\end{equation*}
Substituting the success probability and the conditional mean gives
\begin{equation*}
\begin{aligned}
\underset{\mathbf{t}}{\operatorname{maximize}}
\quad&
\left(
1-\prod_{j=1}^{k}\FRR(t_j)
\right)
\prod_{j=1}^{k}\left(1-\FAR(t_j)\right)
\\
\text{subject to}
\quad&
\frac{
\displaystyle
\sum_{j=1}^{k}
j\left(\prod_{i=1}^{j-1}\FRR(t_i)\right)
\left(1-\FRR(t_j)\right)
}{
\displaystyle
1-\prod_{j=1}^{k}\FRR(t_j)
}
\leq \mu'.
\end{aligned}
\end{equation*}
We restrict the optimization to threshold vectors for which
$P_U(\mathbf{t})>0$, so that the conditional mean is defined.
This optimization maximizes the mechanism’s success probability, while ensuring that successful user authentication sessions require at most~$\mu'$ attempts on average.

Here the constraint plays a different role. Without the constraint, the optimization balances both the~\FRR and~\FAR, leading to a non trivial solution as described in~$(\S\ref{sec:equal_k_attempts})$. The constraint~$\mu'$ caps how long a successful user waits. Unlike the previous problem, this constraint may not bind at all: if the optimal threshold vector already lets users authenticate within~$\mu'$ attempts on average, nothing changes. When it does bind, the optimization enforces thresholds more permissive than the unconstrained problem, leading to a reduction in the overall success probability. Thus, this optimization fixes a usability budget and maximizes the mechanism's success probability within it.

\subsection{Spoofing Attack With Multiple Attempts}\label{sec:mul_attempts_spoofing}
\noindent
We simulate a spoofing attack scenario in which a principal has multiple attempts to authenticate. We use the same genuine and impostor distributions calculated in the spoofing attack scenario using the LivDet2015 database. 

We follow the theoretical analysis and calculate the success probability under multiple attempts of a mechanism that uses a biometric credential only~(Equation~\ref{eqn:success_multiple_attempts}). We then simulate several authentication attempts by a principal~$P$. For each number of attempts~$k$, we calculate the mechanism's success probability as the probability that the user is accepted at least once, while the attacker is rejected on all~$k$ attempts. As the success probability depends on the number of attempts~$k$, we compare the mechanism's success probability under three thresholds selections baselines: (1)~the~\EER threshold reused for all~$k$ attempts, (2)~the single-attempt optimal threshold reused for all~$k$ attempts, and (3)~threshold optimization over all~$k$ attempts, as shown in Figure~\ref{fig:Success_probability_by_k}. The first two baselines represent biometric authentication systems that use a fixed threshold across multiple authentication attempts, either the~\EER threshold or the single-attempt optimal threshold, without optimizing it for multiple attempts.

We find that the success probability of this mechanism, which uses a biometric credential only, increases with the number of attempts~$k$ and converges to 1 when optimizing the threshold over all~$k$ attempts~(Figure~\ref{fig:unconstrained_success_convergence}). The optimization suggests that the threshold should vary with the total number of attempts~$k$. Furthermore, our findings reveal that reusing a single threshold in all~$k$ attempts, either the~\EER threshold or the single-attempt optimal threshold, leads to a reduction of up to 75\% in the mechanism's failure probability. This reduction indicates that the attacker is more likely to succeed in spoofing the biometric credential as the number of attempts increases.

Previously, we noted that for each number of attempts~$k$, the optimal threshold varies with~$k$. Thus, we investigate how the optimal threshold selections are affected when we constrain the mean number of attempts, seeking to maximize the mechanism's success probability. We find that the optimal threshold increases with the number of attempts~$k$, and that the mechanism's success probability increases as the mean number of attempts increases. Nevertheless, system designers seek to minimize the mean number of a user's authentication attempts. 

We explore the tradeoff between the mean number of attempts and the mechanism's success probability. First, we identify the optimal threshold selections that minimize the mean number of user authentication attempts, while requiring the mechanism's success probability to exceed~$P_{\mathrm{success}}^{\min}$. This analysis shows that the greater the required success probability~$P_{\mathrm{success}}^{\min}$, the greater the mean number of user authentication attempts~$\mathbb{E}[N_U]$ required to achieve it. Moreover, as presented in the theoretical analysis, the mechanism's success probability is equal to the success constraint. In addition, our findings reveal that a higher success probability constraint leads to a greater cumulative user success probability and a lower cumulative attacker success probability. This behavior indicates that the selected thresholds tend to be stricter in the initial attempts, with smaller differences between successive thresholds as the attempt index increases. In this way, the mechanism limits the probability of a successful spoofing attack, while relying on the user's cumulative success probability to increase with the number of attempts. These findings are summarized in Figure~\ref{fig:success_constraints_comparison}.

We also analyze the reverse direction, in which a system designer seeks to maximize the achievable security under a usability constraint. Specifically, we constrain the mean number of user authentication attempts~$\mathbb{E}[N_U]$, and seek to maximize the mechanism's success probability. We find that allowing a greater mean number of user authentication attempts~$\mathbb{E}[N_U]$, leads to a lower cumulative success probability for the attacker. 
A higher limit on the mean number of attempts leads to stricter threshold selections in the initial authentication attempts. In this way, the mechanism prioritizes preventing the attacker from spoofing the biometric credential, over accepting the user in the initial attempts. Thus, the user's cumulative authentication probability is initially lower than under a lower limit on the mean number of attempts. However, it increases significantly over subsequent attempts, as the user has more opportunities to obtain a score above the threshold. These findings are shown in Figure~\ref{fig:mean_limits_comparison}. 

\begin{figure}[t]
\centering

% Upper subfigure: probabilities versus maximum attempt count
\begin{subfigure}{\columnwidth}
\centering
\begin{tikzpicture}
\begin{axis}[
    width=0.95\linewidth,
    height=0.65\linewidth,
    xlabel={Number of attempts, $k$},
    ylabel={Probability},
    xmin=1,
    xmax=10,
    ymin=0,
    ymax=1,
    xtick=data,
    grid=major,
    line width=1pt,
    mark size=1.6pt,
    legend style={
        at={(0.5,-0.28)},
        anchor=north,
        legend columns=3,
        draw=black,
        fill=white,
        font=\scriptsize
    }
]

\addplot[green!60!black, solid, mark=*]
table[x=k, y=P_success_max]
{figs/fig_spoofed_users/livdet_fixed_prefix_far_frr_success_by_k.txt};
\addlegendentry{$P_{\mathrm{success}}^{\max}$}

\addplot[blue, dashed, mark=square*]
table[x=k, y=P_success_EER]
{figs/fig_spoofed_users/livdet_fixed_prefix_far_frr_success_by_k.txt};
\addlegendentry{$P_{\mathrm{success}}(t_{\mathrm{EER}})$}

\addplot[red, densely dotted, mark=triangle*]
table[x=k, y=P_success_single_attempt_optimum]
{figs/fig_spoofed_users/livdet_fixed_prefix_far_frr_success_by_k.txt};
\addlegendentry{$P_{\mathrm{success}}(t_{\mathrm{opt}}^{(1)})$}

\addplot[orange, solid, mark=diamond*]
table[x=k, y=product_FRR]
{figs/fig_spoofed_users/livdet_fixed_prefix_far_frr_success_by_k.txt};
\addlegendentry{$\prod_{j=1}^{k}\mathrm{FRR}_j$}

\addplot[purple, solid, mark=pentagon*]
table[x=k, y=product_1_minus_FAR]
{figs/fig_spoofed_users/livdet_fixed_prefix_far_frr_success_by_k.txt};
\addlegendentry{$\prod_{j=1}^{k}(1-\mathrm{FAR}_j)$}

\end{axis}
\end{tikzpicture}
\caption{Cumulative success and rejection probabilities.}
\label{fig:Success_probability_by_k_probabilities}
\end{subfigure}

\medskip

% Lower subfigure: final applied threshold for each k-attempt policy
\begin{subfigure}{\columnwidth}
\centering
\begin{tikzpicture}
\begin{axis}[
    width=0.95\linewidth,
    height=0.50\linewidth,
    xlabel={Number of attempts, $k$},
    ylabel={Selected threshold, $t_k^{(k)}$},
    xmin=1,
    xmax=10,
    xtick=data,
    grid=major,
    line width=1pt,
    mark size=1.6pt
]

\addplot[green!60!black, solid, mark=*]
table[x=k, y=Threshold]
{figs/fig_spoofed_users/livdet_fixed_prefix_far_frr_success_by_k.txt};

\end{axis}
\end{tikzpicture}
\caption{Threshold applied at the final attempt of each policy.}
\label{fig:Success_probability_by_k_thresholds}
\end{subfigure}

\caption{Success probabilities and selected thresholds as a function
of the maximum number of authentication attempts~$k$.
(a)~The optimized cumulative success probability, compared with
reusing the \EER threshold and the single-attempt optimal threshold,
together with the products of user-rejection and attacker-rejection
probabilities.
(b)~The threshold applied at attempt~$k$ of the corresponding
$k$-attempt fixed-prefix policy.
Each value of~$k$ represents a separately optimized policy.}
\label{fig:Success_probability_by_k}
\end{figure}

\pgfplotsset{compat=1.18}

% Document
\begin{figure}[t]
\centering
\begin{tikzpicture}
\begin{axis}[
    width=0.95\columnwidth,
    height=0.65\columnwidth,
    xlabel={Number of attempts, $k$},
    ylabel={Overall success probability},
    xmin=1, xmax=200,
    ymin=0, ymax=1.02,
    xtick={1,25,50,75,100,125,150,175,200},
    ytick={0,0.2,0.4,0.6,0.8,1},
    grid=major,
    legend style={
        at={(0.5,-0.31)},
        anchor=north,
        legend columns=1,
        draw=black,
        fill=white,
        font=\scriptsize
    },
    legend cell align=left
]

\addplot[
    blue,
    thick,
    mark=o,
    mark size=1.8pt,
    mark repeat=10,
    mark options={solid}
]
table[
    x=k,
    y=P_success_optimized
] {figs/fig_spoofed_users/livdet_unconstrained_success_convergence.txt};
\addlegendentry{Optimized overall success}

\end{axis}
\end{tikzpicture}
\caption{Optimized mechanism success probability as a function of
the number of authentication attempts~$k$.
The threshold policy is optimized separately for each~$k$,
without constraints on the conditional mean number of attempts
or the success probability. The dashed line marks probability one.}
\label{fig:unconstrained_success_convergence}
\end{figure}

\pgfplotsset{compat=1.18}

\pgfplotsset{
    constraint axis/.style={
        width=0.9\linewidth,
        height=0.55\linewidth,
        xlabel={Attempt},
        xmin=1, xmax=10,
        xtick={1,...,10},
        grid=major,
        mark size=2pt,
        legend style={
            at={(0.5,-0.31)},
            anchor=north,
            draw=black,
            fill=white,
            font=\scriptsize
        },
        legend cell align=left
    },
    player axis/.style={
        constraint axis,
        ylabel={Cumulative success probability},
        ymin=0, ymax=1,
        legend columns=3,
        every axis plot/.append style={forget plot}
    }
}

% Two-row legend:
% Row 1: three constraint values and their line styles.
% Row 2: user and attacker markers.
\newcommand{\TwoRowConstraintLegend}[3]{%
    \addlegendimage{black, thick, solid, mark=none}
    \addlegendentry{#1}
    \addlegendimage{black, thick, densely dotted, mark=none}
    \addlegendentry{#2}
    \addlegendimage{black, thick, dashed, mark=none}
    \addlegendentry{#3}
    \addlegendimage{blue, only marks, mark=o, mark size=2pt}
    \addlegendentry{User}
    \addlegendimage{red, only marks, mark=square, mark size=2pt}
    \addlegendentry{Attacker}
    \addlegendimage{empty legend}
    \addlegendentry{}
}

% =========================================================
% Figure 1: minimize mean attempts under success constraints
% =========================================================

\begin{figure}[t]
\centering
\def\PolicyData{figs/fig_spoofed_users/livdet_success_constraints_comparison.txt}

% Upper: threshold versus attempt
\begin{subfigure}{\linewidth}
\centering
\begin{tikzpicture}
\begin{axis}[
    constraint axis,
    ylabel={Threshold},
    ymin=0, ymax=100
]

\addplot[black, thick, solid, mark=*, mark size=1pt, mark options={solid}]
table[x=Attempt, y=Threshold_0_6] {\PolicyData};

\addplot[black, thick, densely dotted, mark=*, mark size=1pt, mark options={solid}]
table[x=Attempt, y=Threshold_0_7] {\PolicyData};

\addplot[black, thick, dashed, mark=*, mark size=1pt, mark options={solid}]
table[x=Attempt, y=Threshold_0_75] {\PolicyData};

\end{axis}
\end{tikzpicture}
\caption{Selected thresholds.}
\label{fig:success_constraints_thresholds}
\end{subfigure}

\par\medskip

% Lower: user and attacker success versus attempt
\begin{subfigure}{\linewidth}
\centering
\begin{tikzpicture}
\begin{axis}[player axis]

% Lowest constraint: solid
\addplot[blue, thick, solid, mark=o, mark options={solid}]
table[x=Attempt, y=P_user_success_0_6] {\PolicyData};

\addplot[red, thick, solid, mark=square, mark options={solid}]
table[x=Attempt, y=P_attacker_success_0_6] {\PolicyData};

% Middle constraint: dotted
\addplot[blue, thick, densely dotted, mark=o, mark options={solid}]
table[x=Attempt, y=P_user_success_0_7] {\PolicyData};

\addplot[red, thick, densely dotted, mark=square, mark options={solid}]
table[x=Attempt, y=P_attacker_success_0_7] {\PolicyData};

% Highest constraint: dashed
\addplot[blue, thick, dashed, mark=o, mark options={solid}]
table[x=Attempt, y=P_user_success_0_75] {\PolicyData};

\addplot[red, thick, dashed, mark=square, mark options={solid}]
table[x=Attempt, y=P_attacker_success_0_75] {\PolicyData};

\TwoRowConstraintLegend
    {$P'=0.6$}
    {$P'=0.7$}
    {$P'=0.75$}

\end{axis}
\end{tikzpicture}
\caption{Cumulative user and attacker acceptance probabilities.}
\label{fig:success_constraints_probabilities}
\end{subfigure}

\caption{Policies minimizing the conditional mean number of user
authentication attempts subject to
$P_{\mathrm{success}}\geq P'$, with $k=10$.
The upper panel shows selected thresholds; the lower panel shows
the probability of at least one user or attacker acceptance through
each attempt. Line styles identify constraint values, while blue
circles represent the user and red squares represent the attacker.
The constraint applies to the complete policy.}
\label{fig:success_constraints_comparison}
\end{figure}

% =========================================================
% Figure 2: maximize success under mean-attempt constraints
% =========================================================

\begin{figure}[t]
\centering
\def\PolicyData{figs/fig_spoofed_users/livdet_mean_limits_comparison.txt}

% Upper: threshold versus attempt
\begin{subfigure}{\linewidth}
\centering
\begin{tikzpicture}
\begin{axis}[
    constraint axis,
    ylabel={Threshold},
    ymin=0, ymax=100
]

\addplot[black, thick, solid, mark=*, mark size=1pt, mark options={solid}]
table[x=Attempt, y=Threshold_1_5] {\PolicyData};

\addplot[black, thick, densely dotted, mark=*, mark size=1pt, mark options={solid}]
table[x=Attempt, y=Threshold_1_9] {\PolicyData};

\addplot[black, thick, dashed, mark=*, mark size=1pt, mark options={solid}]
table[x=Attempt, y=Threshold_2_4] {\PolicyData};

\end{axis}
\end{tikzpicture}
\caption{Selected thresholds.}
\label{fig:mean_limits_thresholds}
\end{subfigure}

\par\medskip

% Lower: user and attacker success versus attempt
\begin{subfigure}{\linewidth}
\centering
\begin{tikzpicture}
\begin{axis}[player axis]

% Lowest constraint: solid
\addplot[blue, thick, solid, mark=o, mark options={solid}]
table[x=Attempt, y=P_user_success_1_5] {\PolicyData};

\addplot[red, thick, solid, mark=square, mark options={solid}]
table[x=Attempt, y=P_attacker_success_1_5] {\PolicyData};

% Middle constraint: dotted
\addplot[blue, thick, densely dotted, mark=o, mark options={solid}]
table[x=Attempt, y=P_user_success_1_9] {\PolicyData};

\addplot[red, thick, densely dotted, mark=square, mark options={solid}]
table[x=Attempt, y=P_attacker_success_1_9] {\PolicyData};

% Highest constraint: dashed
\addplot[blue, thick, dashed, mark=o, mark options={solid}]
table[x=Attempt, y=P_user_success_2_4] {\PolicyData};

\addplot[red, thick, dashed, mark=square, mark options={solid}]
table[x=Attempt, y=P_attacker_success_2_4] {\PolicyData};

\TwoRowConstraintLegend
    {$\mu'=1.5$}
    {$\mu'=1.9$}
    {$\mu'=2.4$}

\end{axis}
\end{tikzpicture}
\caption{Cumulative user and attacker acceptance probabilities.}
\label{fig:mean_limits_probabilities}
\end{subfigure}

\caption{Policies maximizing the mechanism's success probability
subject to $\mathbb{E}[N_U\mid N_U\leq k]\leq\mu'$,
with $k=10$.
The upper panel shows selected thresholds; the lower panel shows
the probability of at least one user or attacker acceptance through
each attempt. Line styles identify constraint values, while blue
circles represent the user and red squares represent the attacker.
The constraint applies to the complete policy.}
\label{fig:mean_limits_comparison}
\end{figure}


% \section{Empirical Evaluation}\label{sec:empirical_evaluation}
% \noindent
% We now empirically evaluate different mechanisms that use an atomic credential and a biometric credential under two attack scenarios. In each scenario, we compare the mechanism's success probability when the threshold is set to the~\EER and when it is optimized. We first simulate a baseline attack in which an attacker attempts to authenticate using their own fingerprints~$(\S\ref{sec:baseline_attack})$. We show that although optimizing the threshold increases the mechanism's success probability, the improvement is limited. We then present a more sophisticated attack in which an attacker creates a fake fingerprint using spoofing techniques~$(\S\ref{sec:spoofing_attack})$ and show that the improvement becomes substantially larger. To reflect real-world authentication systems, we evaluate spoofing attack when multiple attempts are allowed~$(\S\ref{sec:mul_attempts_spoofing})$. We show that in such cases, the threshold should be optimized with respect to the maximum number of allowed attempts for siginificant improvement in the mechanism's success probability.

% \subsection{Baseline Attack Scenario}\label{sec:baseline_attack}
% \noindent
% To simulate a biometric credential, we use the public real-world~\textit{LivDet2015} dataset of fingerprints to compare fingerprint verification algorithms~\cite{LivDet2015,LivDet2013,LivDet2011,LivDet2009}. This dataset contains multiple fingerprint scans of each user, acquired using different types of biometric sensors. We divide each user's scans in the dataset into a reference set, simulating the reference scans stored in advance, and a candidate set, simulating the scans obtained by the sensor. To score the fingerprint scan images, we use the standard fingerprint verification algorithm~\textit{SourceAFIS}~\cite{sourceafis}. This algorithm compares the relative positions and orientations of fingerprint ridge endings and bifurcations across two images. We evaluate all users who appear in~$\textit{LivDet2015}$ dataset whose fingerprints obtained by the same sensor. For each user, we randomly select five fingerprint scan images for the reference set, while the candidate set comprises the remaining images in the dataset. We compute a similarity score for each candidate scan by running the algorithm against each of the five reference images, and choose the maximum score as the candidate's final score~\cite{hong1998,Cappelli2006,max_score2009}. The genuine comparisons are all comparisons between the user's candidate images and the user's reference images. By contrast, the imposter comparisons are all comparisons between the users' reference images and the scans of the other users in the database.

% Calculating the mechanism's success probability requires the biometric credential's fault probabilities, which are functions of the~\FAR and~\FRR. We therefore follow the evaluation procedure used by Jain et al.~\cite{Anil1997}, Cappelli et al.~\cite{Cappelli2006} and Slobodan et al.~\cite{Slob2005}. We produce histograms of the genuine and imposter comparisons, as shown in Figure~\ref{fig:sourceafis_histogram}. Then, we calculate the~\FAR and~\FRR as functions of the threshold~$t$, as shown in Figure~\ref{fig:sourceafis_far_frr}.

% Before comparing different mechanisms, we validate that the algorithm's performance is representative of the state-of-the-art fingerprint verification algorithms. The algorithm we use achieves an~\EER of approximately~3.8\%, similar to the best results reported in the~\textit{LivDet2015}, and to those achieved by fingerprint systems used in industry and academia when evaluating different users~\cite{Maio2002,Cappelli2006}. We calculate the biometric credential fault probabilities~(Equation~\ref{eqn:safe_prob}) and the success probability of a mechanism that uses a biometric credential only~(Equation~\ref{eqn:safesuccess}). We find that the optimal threshold differs from the~\EER point and achieves a higher success probability, as shown in Figure~\ref{fig:bio_only_sourceafis_reg_scoring}. 

% \newcommand{\sourceafisplotspanel}{%
% \begin{figure*}[!t]
% \centering
% \begin{subfigure}[t]{0.48\textwidth}
% \centering
% \begin{tikzpicture}
% \begin{axis}[
%     ybar,
%     legend to name=histogramlegend,
%     bar width=3pt,
%     width=0.86\linewidth,
%     height=0.40\linewidth,
%     xlabel={Normalized similarity score (\%)},
%     ylabel={Impostor count},
%     ylabel style={red},
%     grid=major,
%     xmin=0,
%     xmax=100,
%     ymin=0,
%     ymax=9000,
%     enlarge x limits=false,
%     enlarge y limits=false,
%     tick style={draw=none},
%     clip=true
% ]

% \addplot+[ybar, fill=red, draw=black]
% table [x=bin, y=impostor] {figs/fig_different_users_sourceafsi/sourceafis_histogram_data.txt};

% \end{axis}

% \begin{axis}[
%     ybar,
%     bar width=3pt,
%     legend to name=histogramlegend,
%     legend columns=2,
%     width=0.86\linewidth,
%     height=0.40\linewidth,
%     xlabel={Normalized similarity score (\%)},
%     yticklabel pos=right,
%     ylabel={Genuine count},
%     ylabel style={blue},
%     grid=none,
%     xmin=0,
%     xmax=100,
%     ymin=0,
%     ymax=80,
%     enlarge x limits=false,
%     enlarge y limits=false,
%     tick style={draw=none},
%     clip=true
% ]

% \addplot+[ybar, fill=blue, fill opacity=0.5]
% table [x=bin, y=genuine] {figs/fig_different_users_sourceafsi/sourceafis_histogram_data.txt};
% \addlegendentry{Genuine}

% \addlegendimage{fill=red, mark=none}
% \addlegendentry{Impostor}
% \end{axis}
% \end{tikzpicture}
% \medskip
% \pgfplotslegendfromname{histogramlegend}
% \caption{Score histogram for genuine and impostor comparisons.}
% \label{fig:sourceafis_histogram}
% \end{subfigure}
% \hfill
% \begin{subfigure}[t]{0.48\textwidth}
% \centering
% \begin{tikzpicture}
% \begin{axis}[
%     scale only axis,
%     width=0.86\linewidth,
%     height=0.40\linewidth,
%     xlabel={Threshold (\%)},
%     ylabel={Error Rate},
%     ymin=0,
%     ymax=1,
%     grid=major,
%     legend style={
%         at={(0.5,-0.30)},
%         anchor=north,
%         legend columns=2,
%         draw=black,
%         fill=white,
%         font=\scriptsize
%     }
% ]

% \pgfplotsset{
%     impostor curve/.style={
%         thick, dashed
%     },
%     genuine curve/.style={
%         thick, solid
%     },
%     uniform model/.style={purple},
%     quadratic model/.style={orange},
%     normal model/.style={cyan}
% }

% % Empirical FAR: discrete histogram integral
% \addplot[
%     only marks,
%     red,
%     mark=*,
%     mark size=1.3pt
% ]
% table[x=T, y=FAR] {figs/fig_different_users_sourceafsi/sourceafis_far_frr_data.txt};
% \addlegendentry{Empirical \FAR}

% % Empirical FRR: discrete histogram integral
% \addplot[
%     only marks,
%     blue,
%     mark=*,
%     mark size=1.3pt
% ]
% table[x=T, y=FRR] {figs/fig_different_users_sourceafsi/sourceafis_far_frr_data.txt};
% \addlegendentry{Empirical \FRR}

% % Uniform: squares
% \addplot[impostor curve, uniform model]
% table[x=T, y=FAR_uniform]
% {figs/fig_different_users_sourceafsi/sourceafis_far_frr_data.txt};
% \addlegendentry{Uniform \FAR}

% \addplot[genuine curve, uniform model]
% table[x=T, y=FRR_uniform]
% {figs/fig_different_users_sourceafsi/sourceafis_far_frr_data.txt};
% \addlegendentry{Uniform \FRR}

% % Quadratic: diamonds
% \addplot[impostor curve, quadratic model]
% table[x=T, y=FAR_quadratic]
% {figs/fig_different_users_sourceafsi/sourceafis_far_frr_data.txt};
% \addlegendentry{Quadratic \FAR}

% \addplot[genuine curve, quadratic model]
% table[x=T, y=FRR_quadratic]
% {figs/fig_different_users_sourceafsi/sourceafis_far_frr_data.txt};
% \addlegendentry{Quadratic \FRR}

% % Normal: triangles
% \addplot[impostor curve, normal model]
% table[x=T, y=FAR_normal]
% {figs/fig_different_users_sourceafsi/sourceafis_far_frr_data.txt};
% \addlegendentry{Normal \FAR}

% \addplot[genuine curve, normal model]
% table[x=T, y=FRR_normal]
% {figs/fig_different_users_sourceafsi/sourceafis_far_frr_data.txt};
% \addlegendentry{Normal \FRR}

% % Empirical EER: black star
% \addplot[
%     black,
%     only marks,
%     mark=star,
%     mark size=3.5pt,
%     mark options={line width=2pt}
% ]
% table[x=T_eer, y=EER]
% {figs/fig_different_users_sourceafsi/sourceafis_far_frr_points.txt};
% \addlegendentry{\EER}

% \end{axis}
% \end{tikzpicture}
% \caption{Empirical \FAR and \FRR as functions of the decision threshold, compared with the uniform and quadratic probabilistic distribution functions obtained from the KS fits.}
% \label{fig:sourceafis_far_frr}
% \end{subfigure}

% \par\medskip

% \pgfplotsset{
%     sourceafissuccess/.style={
%         scale only axis,
%         width=0.82\linewidth,
%         height=0.36\linewidth,
%         xlabel={Threshold (\%)},
%         ylabel={Success Probability},
%         xmin=0,
%         xmax=100,
%         ymin=0,
%         ymax=1,
%         grid=major,
%         legend style={
%             at={(0.5,-0.3)},
%             anchor=north,
%             legend columns=3,
%             draw=black,
%             fill=white,
%             font=\scriptsize
%         },
%         filter discard warning=false
%     }
% }

% \begin{subfigure}[t]{0.48\textwidth}
% \centering
% \begin{tikzpicture}
% \begin{axis}[
%     sourceafissuccess
% ]

% \addplot[thick, blue, mark=none]
% table [x=T, y=P_success] {figs/fig_different_users_sourceafsi/sourceafis_p_success_data.txt};
% \addlegendentry{$P_{\mathrm{success}}(t)$}

% \addplot[only marks, mark=square*, mark size=1.4pt, myGreen]
% table [x=T_eer, y=P_eer] {figs/fig_different_users_sourceafsi/sourceafis_p_success_points.txt};
% \addlegendentry{$P_{\mathrm{success}}(t_{\EER})$}

% \addplot[only marks, mark=*, mark size=1.6pt, myRed]
% table [x=T_opt, y=P_opt] {figs/fig_different_users_sourceafsi/sourceafis_p_success_points.txt};
% \addlegendentry{Optimal point}

% \end{axis}
% \end{tikzpicture}
% \caption{Single biometric credential mechanism's success.}
% \label{fig:bio_only_sourceafis_reg_scoring}
% \end{subfigure}
% \hfill
% \begin{subfigure}[t]{0.48\textwidth}
% \centering
% \begin{tikzpicture}
% \begin{axis}[
%     sourceafissuccess
% ]

% \addplot[thick, blue, mark=none]
% table [x=T, y=P_and] {figs/fig_different_users_sourceafsi/sourceafis_success_and_or_data.txt};
% \addlegendentry{$P_{\mathrm{success}}^{\mathrm{AND}}(t)$}

% \addplot[thick, myOrange, mark=none]
% table [x=T, y=P_or] {figs/fig_different_users_sourceafsi/sourceafis_success_and_or_data.txt};
% \addlegendentry{$P_{\mathrm{success}}^{\mathrm{OR}}(t)$}

% \addplot[only marks, mark=*, mark size=1.5pt, blue]
% table [x=T_and_opt, y=P_and_opt] {figs/fig_different_users_sourceafsi/sourceafis_success_and_or_points.txt};
% \addlegendentry{AND optimum}

% \addplot[only marks, mark=x, mark size=1.9pt, myOrange]
% table [x=T_or_opt, y=P_or_opt] {figs/fig_different_users_sourceafsi/sourceafis_success_and_or_points.txt};
% \addlegendentry{OR optimum}

% \addplot[only marks, mark=square*, mark size=1.4pt, myGreen]
% table [x=T_eer, y=P_and_eer] {figs/fig_different_users_sourceafsi/sourceafis_success_and_or_points.txt};
% \addlegendentry{AND at \EER}

% \addplot[only marks, mark=diamond*, mark size=1.5pt, myRed]
% table [x=T_eer, y=P_or_eer] {figs/fig_different_users_sourceafsi/sourceafis_success_and_or_points.txt};
% \addlegendentry{OR at \EER}

% \addplot[
%     thick,
%     black,
%     dashed,
%     domain=0:100,
% ] {0.75};
% \addlegendentry{Atomic credential only}

% \end{axis}
% \end{tikzpicture}
% \caption{biometric credential integrated with an atomic credential into AND and OR mechanisms.}
% \label{fig:bio_atm_sourceafis_reg_scoring}
% \end{subfigure}

% % \Description{Success probability curves as a function of threshold, including \EER and optimal operating points.}
% \caption{
% Success probability as a function of the threshold.
% The AND and OR mechanisms achieve success probabilities of 
% $0.78$ and $0.82$ at the \EER point, and 
% $0.81$ and $0.84$ at their respective optimal operating points.
% }
% \label{fig:sourceafis_success}
% \end{figure*}
% }

% We next evaluate mechanisms that use an atomic credential and a biometric credential. We simulate an atomic credential with arbitrary fault probabilities~$P^{1}_\text{safe} = 0.75$, $P^{1}_\text{leak} = 0.1$, $P^{1}_\text{loss} = 0.1$, and $P^{1}_\text{stolen} = 0.05$. Then, we calculate the success probability of the \emph{AND mechanism}~(Equation~\ref{eqn:andsuccess}), and that of the \emph{OR mechanism}~(Equation~\ref{eqn:orsuccess}). We illustrate the success probability functions of mechanisms, that use an atomic credential and a biometric credential as functions of~$t$ in Figure~\ref{fig:bio_atm_sourceafis_reg_scoring}. We show that the optimal threshold for each of the mechanisms differs from the~\EER point. Moreover, both mechanisms operating at their optimal thresholds achieve success probabilities that are~8-9\% higher than that of the atomic credential only mechanism, demonstrating the benefit of using the biometric credential in addition. 

% However, the success probabilities of all evaluated mechanisms, when operated at their optimal thresholds, are not significantly greater than those of the corresponding mechanisms operating at the~\EER point. This can be explained by the fact that the~\EER of the biometric credential is already very low, indicating high separability between a user and an attacker and thus, leaving small room for improvement. In this scenario, we tested different users against each other. However, our main objective is to increase the mechanism's success probability in scenarios where an imposter attempts to spoof the user's templates, and authenticate to the user's account. Hence, we expect the imposter to use spoofing tools, producing samples whose score distribution is much more similar to the genuine distribution, which might lead to an increased~\EER point. 

% \begingroup
% \makeatletter
% \setlength{\@dblfptop}{0.5in}
% \makeatother
% \sourceafisplotspanel
% \clearpage
% \endgroup

% \newcommand{\livdetplotblock}{%
% \begin{strip}
% \vspace{-0.6\baselineskip}
% \centering
% \begin{minipage}{\textwidth}
% \captionsetup{type=figure}
% \centering
% \begin{subfigure}[t]{0.48\textwidth}
% \centering
% \begin{tikzpicture}
% % Left axis: impostor counts
% \begin{axis}[
%     ybar,
%     legend to name=livdethistogramlegend,
%     bar width=2pt,
%     width=0.86\linewidth,
%     height=0.45\linewidth,
%     xlabel={Normalized similarity score (\%)},
%     ylabel={Impostor count},
%     ylabel style={red},
%     grid=major,
%     xmin=0,
%     xmax=100,
%     ymin=0,
%     ymax=600,
%     enlarge x limits=false,
%     tick style={draw=none},
%     clip=true
% ]
% \addplot+[ybar,fill=red,draw=black]
% table[x=score,y=impostor] 
% {figs/fig_spoofed_users/livdet_sourceafis_histogram_best_of_probes_data.txt};
% \end{axis}
% % Right axis: genuine counts
% \begin{axis}[
%     ybar,
%     legend to name=livdethistogramlegend,
%     legend columns=2,
%     bar width=2pt,
%     width=0.86\linewidth,
%     height=0.45\linewidth,
%     xlabel={Normalized similarity score (\%)},
%     ylabel={Genuine count},
%     ylabel style={blue},
%     yticklabel pos=right,
%     grid=none,
%     xmin=0,
%     xmax=100,
%     ymin=0,
%     ymax=80,
%     enlarge x limits=false,
%     tick style={draw=none},
%     clip=true
% ]
% \addplot+[ybar,fill=blue,draw=blue,fill opacity=0.5]
% table[x=score,y=genuine] 
% {figs/fig_spoofed_users/livdet_sourceafis_histogram_best_of_probes_data.txt};
% \addlegendentry{Genuine}
% \addlegendimage{fill=red,draw=black,mark=none}
% \addlegendentry{Impostor}
% \end{axis}
% \end{tikzpicture}
% \smallskip
% \pgfplotslegendfromname{livdethistogramlegend}
% \caption{LivDet score histogram using the best score over probe templates.}
% \label{fig:livdet_sourceafis_histogram}
% \end{subfigure}\hfill%
% \begin{subfigure}[t]{0.48\textwidth}
% \centering
% \begin{tikzpicture}
% \begin{axis}[
%     width=0.86\linewidth,
%     height=0.45\linewidth,
%     xlabel={Similarity score},
%     ylabel={Probability density},
%     xmin=0,
%     xmax=100,
%     ymin=0,
%     grid=major,
%     legend style={font=\small},
% ]

% % Genuine: fitted normal PDF
% \addplot[blue, thick, domain=0:100, samples=300]
% {1/(17.7711*sqrt(2*pi))*exp(-0.5*((x-30.975)/17.7711)^2)};
% \addlegendentry{Genuine: normal}

% % Impostor: fitted normal PDF
% \addplot[red, thick, domain=0:100, samples=300]
% {1/(11.5956*sqrt(2*pi))*exp(-0.5*((x-10.37179232)/11.5956)^2)};
% \addlegendentry{Impostor: normal}

% \end{axis}
% \end{tikzpicture}
% \caption{Normal probability-density functions fitted to the genuine and impostor scores using the KS estimation.}
% \label{fig:livdet_ks_estimation}
% \end{subfigure}



% \begin{subfigure}[t]{0.48\textwidth}
% \centering
% \begin{tikzpicture}
% \begin{axis}[
%     scale only axis,
%     width=0.86\linewidth,
%     height=0.4\linewidth,
%     xlabel={Threshold (\%)},
%     ylabel={Error Rate},
%     xmin=0,
%     xmax=100,
%     ymin=0,
%     ymax=1,
%     grid=major,
%     legend style={
%         at={(0.5,-0.45)},
%         anchor=north,
%         legend columns=2,
%         draw=black,
%         fill=white,
%         font=\scriptsize
%     }
% ]

% \pgfplotsset{
%     impostor curve/.style={
%         thick, dashed
%     },
%     genuine curve/.style={
%         thick, solid
%     },
%     uniform model/.style={purple},
%     quadratic model/.style={orange},
%     normal model/.style={cyan}
% }

% % Empirical FAR: discrete histogram integral
% \addplot[
%     only marks,
%     red,
%     mark=*,
%     mark size=1.3pt
% ]
% table[x=T, y=FAR] {figs/fig_spoofed_users/livdet_sourceafis_far_frr_best_of_probes_data.txt};
% \addlegendentry{Empirical \FAR}

% % Empirical FRR: discrete histogram integral
% \addplot[
%     only marks,
%     blue,
%     mark=*,
%     mark size=1.3pt
% ]
% table[x=T, y=FRR] {figs/fig_spoofed_users/livdet_sourceafis_far_frr_best_of_probes_data.txt};
% \addlegendentry{Empirical \FRR}

% % Uniform: squares
% \addplot[impostor curve, uniform model]
% table[x=T, y=FAR_uniform]
% {figs/fig_spoofed_users/livdet_sourceafis_far_frr_best_of_probes_data.txt};
% \addlegendentry{Uniform \FAR}

% \addplot[genuine curve, uniform model]
% table[x=T, y=FRR_uniform]
% {figs/fig_spoofed_users/livdet_sourceafis_far_frr_best_of_probes_data.txt};
% \addlegendentry{Uniform \FRR}

% % Quadratic: diamonds
% \addplot[impostor curve, quadratic model]
% table[x=T, y=FAR_quadratic]
% {figs/fig_spoofed_users/livdet_sourceafis_far_frr_best_of_probes_data.txt};
% \addlegendentry{Quadratic \FAR}

% \addplot[genuine curve, quadratic model]
% table[x=T, y=FRR_quadratic]
% {figs/fig_spoofed_users/livdet_sourceafis_far_frr_best_of_probes_data.txt};
% \addlegendentry{Quadratic \FRR}

% % Normal: triangles
% \addplot[impostor curve, normal model]
% table[x=T, y=FAR_normal]
% {figs/fig_spoofed_users/livdet_sourceafis_far_frr_best_of_probes_data.txt};
% \addlegendentry{Normal \FAR}

% \addplot[genuine curve, normal model]
% table[x=T, y=FRR_normal]
% {figs/fig_spoofed_users/livdet_sourceafis_far_frr_best_of_probes_data.txt};
% \addlegendentry{Normal \FRR}

% % Empirical EER: black star
% \addplot[
%     black,
%     only marks,
%     mark=star,
%     mark size=3.5pt,
%     mark options={line width=2pt}
% ]
% table[x=T_eer, y=EER]
% {figs/fig_spoofed_users/livdet_sourceafis_far_frr_best_of_probes_points.txt};
% \addlegendentry{\EER}

% \end{axis}
% \end{tikzpicture}

% \caption{Empirical \FAR and \FRR as functions of the threshold, compared with the uniform and quadratic PDFs obtained from the KS fits. The empirical \EER point is marked and occurs at $\approx 24.8\%$.}
% \label{fig:livdet_sourceafis_far_frr}
% \end{subfigure}\hfill%
% \begin{subfigure}[t]{0.48\textwidth}
% \centering
% \raisebox{7\baselineskip}[0pt][0pt]{%
% \begin{minipage}[t]{\linewidth}
% \centering

% \pgfplotsset{
%     livdetsuccess/.style={
%         scale only axis,
%         width=0.82\linewidth,
%         height=0.25\linewidth,
%         xlabel={Threshold (\%)},
%         ylabel={Success Probability},
%         xmin=0,
%         xmax=100,
%         ymin=0,
%         ymax=1,
%         grid=major,
%         legend pos=north east,
%         legend style={
%             legend columns=3,
%             at={(0.5,-0.45)},
%             anchor=north,
%             draw=black, 
%             fill=white, 
%             font=\scriptsize},
%         each nth point=2,
%         filter discard warning=false
%     }
% }

% \begin{tikzpicture}
% \begin{axis}[livdetsuccess, xlabel={}]

% \addplot[thick, blue, mark=none]
% table [x=T, y=P_success] {figs/fig_spoofed_users/livdet_sourceafis_p_success_best_of_probes_data.txt};
% \addlegendentry{$P_{\mathrm{success}}(t)$}

% \addplot[only marks, mark=square*, mark size=1.4pt, myGreen]
% table [x=T_eer, y=P_eer] {figs/fig_spoofed_users/livdet_sourceafis_p_success_best_of_probes_points.txt};
% \addlegendentry{$P_{\mathrm{success}}(t_{\EER})$}

% \addplot[only marks, mark=*, mark size=1.6pt, myRed]
% table [x=T_opt, y=P_opt] {figs/fig_spoofed_users/livdet_sourceafis_p_success_best_of_probes_points.txt};
% \addlegendentry{Optimal point}

% \end{axis}
% \end{tikzpicture}
% \par\scriptsize Single biometric credential.\par

% \smallskip

% \begin{tikzpicture}
% \begin{axis}[livdetsuccess]

% \addplot[thick, blue, mark=none]
% table [x=T, y=P_and] {figs/fig_spoofed_users/livdet_sourceafis_success_and_or_best_of_probes_data.txt};
% \addlegendentry{$P_{\mathrm{success}}^{\mathrm{AND}}(t)$}

% \addplot[thick, myOrange, mark=none]
% table [x=T, y=P_or] {figs/fig_spoofed_users/livdet_sourceafis_success_and_or_best_of_probes_data.txt};
% \addlegendentry{$P_{\mathrm{success}}^{\mathrm{OR}}(t)$}

% \addplot[only marks, mark=*, mark size=1.5pt, blue]
% table [x=T_and_opt, y=P_and_opt] {figs/fig_spoofed_users/livdet_sourceafis_success_and_or_best_of_probes_points.txt};
% \addlegendentry{AND optimal point}

% \addplot[only marks, mark=x, mark size=1.9pt, myOrange]
% table [x=T_or_opt, y=P_or_opt] {figs/fig_spoofed_users/livdet_sourceafis_success_and_or_best_of_probes_points.txt};
% \addlegendentry{OR optimal point}

% \addplot[only marks, mark=square*, mark size=1.4pt, myGreen]
% table [x=T_eer, y=P_and_eer] {figs/fig_spoofed_users/livdet_sourceafis_success_and_or_best_of_probes_points.txt};
% \addlegendentry{AND at \EER}

% \addplot[only marks, mark=diamond*, mark size=1.5pt, myRed]
% table [x=T_eer, y=P_or_eer] {figs/fig_spoofed_users/livdet_sourceafis_success_and_or_best_of_probes_points.txt};
% \addlegendentry{OR at \EER}

% \addplot[
%     thick,
%     black,
%     dashed,
%     domain=0:100,
% ] {0.8};
% \addlegendentry{Atomic credential only}

% \end{axis}
% \end{tikzpicture}
% \par\scriptsize AND and OR mechanisms.\par

% % \Description{Success probability as a function of threshold for the LivDet SourceAFIS scores using the best score over probe templates.}
% \caption{
% Success probability as a function of the decision threshold using the best score over probe templates.
% The AND mechanism achieves a maximal success probability of $0.834$, while the OR mechanism achieves $0.816$.
% At the \EER operating point, the success probabilities are $0.713$ and $0.649$ for the AND and OR mechanisms, respectively.
% }
% \label{fig:livdet_sourceafis_success}
% \end{minipage}%
% }%
% \end{subfigure}
% \par\vspace{9\baselineskip}% Reserve space for the raised right-hand panel and its caption.
% \caption{LivDet spoofing-attack evaluation using the best score over probe templates.}
% \label{fig:livdet_evaluation_panel}
% \end{minipage}
% \end{strip}
% }

% \raggedbottom
% \refstepcounter{subsection}
% \subsection*{\thesubsection.\quad Spoofing Attack Scenario}\label{sec:spoofing_attack}
% \makeatletter
% \immediate\write\@auxout{\string\newlabel{sec:spoofing_attack}{{\thesubsection}{\thepage}{Spoofing Attack Scenario}{\@currentHref}{}}}
% \immediate\write\@auxout{\string\@writefile{toc}{\string\contentsline {subsection}{\string\numberline {\thesubsection}Spoofing Attack Scenario}{\thepage}{\@currentHref}\string\protected@file@percent}}
% \makeatother
% \noindent
% To support this claim, we use the LivDet2015 database \cite{LivDet2015,LivDet2013,LivDet2011,LivDet2009}, which contains a large-scale live and spoof fingerprint samples. 
% We follow the same methodology as before and use the $\text{SourceAFIS}$ algorithm~\cite{sourceafis}. As reference images, we randomly select six Digital Persona scanner images of each user in the database, while the probe set of each user comprises the remaining real and fake images. The geniune comparisons are all comparisons between the reference and real images in the probe set.
% For each probe image, we define the similarity score as the maximum score among 6 reference images. 
% We simulate a spoofing attack by using fake fingerprint images of each user produced by~\textit{Gelatine}, \textit{Woodglue}, \textit{Latex} and~\textit{Ecoflex} and match them against the reference images to obtain the imposter score distribution. We build the resulting histogram of a user and an imposter (Figure~\ref{fig:livdet_sourceafis_histogram}) and then calculate the~\FAR-\FRR error rates to compute the~\EER point (Figure~\ref{fig:livdet_sourceafis_far_frr}).
% We assess how closely the theoretical models~$(\S\ref{sec:mechanisms})$ approximate the empirical results by fitting uniform, quadratic, and normal distributions to the genuine and impostor score histograms. For each fitted distribution, we calculate the Kolmogorov-Smirnov~(KS) statistic, which measures the maximum absolute difference between the empirical and fitted cumulative distribution functions~\cite{Massey1951}. The ruslting KS statistics range from $0.04$ to $0.07$ for the genuine scores and are approximately $0.17$ for the impostor scores, providing empirical support for using our theoretical models as approximations of the observed score distributions. 
% As we expect, the~\EER increased significantly to 22.1\% indicating the user and imposter distributions are less separable and more overlapping in spoofing attacks. To understand the implications of a spoofing attack, we calculate the success probability of a mechanism uses a biometric credential and a mechanism uses an atomic credential and a biometric credential, as shown in Figure~\ref{fig:livdet_sourceafis_success}. We find that the optimal operating point reducing the failure probability between 20\% to 72\% compared to the~\EER point when the mechanism uses both an atomic and a biometric credential. Moreover, we find that a mechanism uses an atomic credential and a biometric credential at the~\EER threshold, might achieve lower success probability than that of a mechanism uses an atomic credential only, or a biometric credential only. These results demonstrating that non-optimizing the threshold of the biometric credential, when using multiple credentials in a mechanism, might drawback the mechanism's success probability. These results are summarized in Table~\ref   {tab:psafe_sweep_improvement}.

% \livdetplotblock
% \clearpage

% \begin{strip}
% \centering
% \begin{minipage}{\textwidth}
% \captionsetup{type=table}
%     \centering
%     \caption{Comparison between the success probabilities obtained at the EER threshold and the maximal success probabilities. Bold EER values indicate the worse composition, whereas bold maximal values indicate the best composition. The bold safe column indicating the success probability of a mechanism uses an atomic credential only. The improvement column shows the percentage of improvement between the failure probability at the optimal mechanism success and the failure probability at the maximal mechanism success at EER threshold.}
%     \label{tab:psafe_sweep_improvement}
%     \begin{tabular}{cccccccccc}
%         \toprule
%         $\mathbf{P_{\mathrm{safe}}}$ & $P_{\mathrm{loss}}$ & $P_{\mathrm{leak}}$ & $P_{\mathrm{theft}}$ & AND at EER & OR at EER & Max AND & Max OR & Improvement \\
%         \midrule
%         \textbf{0.60} & 0.30 & 0.05 & 0.05 & \textbf{0.45} & 0.58 & 0.60 & \textbf{0.66} & 20.25\% \\
%         \textbf{0.60} & 0.05 & 0.30 & 0.05 & 0.57 & \textbf{0.45} & \textbf{0.68} & 0.60 & 24.94\% \\
%         \textbf{0.75} & 0.15 & 0.05 & 0.05 & \textbf{0.56} & 0.61 & 0.75 & \textbf{0.77} & 40.39\% \\
%         \textbf{0.82} & 0.05 & 0.08 & 0.05 & 0.62 & \textbf{0.61} & \textbf{0.824} & 0.82 & 53.20\% \\
%         \textbf{0.90} & 0.00 & 0.05 & 0.05 & 0.66 & \textbf{0.64} & 0.89 & \textbf{0.905} & 72.19\% \\
%         \bottomrule
%     \end{tabular}
% \end{minipage}
% \end{strip}
% \flushbottom

% \subsection{Spoofing Attack With Multiple Attempts}\label{sec:mul_attempts_spoofing}
% \noindent
% We simulate a spoofing attack scenario in which a principal has multiple attempts to authenticate. We use the same genuine and impostor distributions calculated in the previous section~$(\S\ref{sec:spoofing_attack})$ using the LivDet2015 database. 

% We follow the theoretical analysis in~$(\S\ref{sec:multiple_auth_attempts})$, and calculate the success probability under multiple attempts of a mechanism that uses a biometric credential only~(Equation~\ref{eqn:success_multiple_attempts}). We then simulate several authentication attempts of a principal~$P$. For each number of attempts~$k$, we calculate the mechanism's success probability as the probability that the user is accepted at least once, while the attacker is rejected on all~$k$ attempts. As the success probability depends on the number of attempts~$k$, we compare the mechanism's success probability under three threshold selections baselines: (1)~the~\EER threshold reused for all~$k$ attempts, (2)~the single attempt optimal threshold reused for all~$k$ attempts, and (3)~threshold optimization over all~$k$ attempts, as shown in Figure~\ref{fig:Success_probability_by_k}. The first two baselines represent biometric authentication systems that use a fixed threshold across multiple authentication attempts, either the~\EER threshold or the single attempt optimal threshold, without optimizing it for multiple attempts.

% We find that the success probability of this mechanism, which uses a biometric credential only, increases with the number of attempts~$k$ and converges to 1, when optimizing the threshold over all~$k$ attempts~(Figure~\ref{fig:unconstrained_success_convergence}). The optimization suggests that the threshold should remain constant across all~$k$ attempts, but vary with the total number of attempts~$k$. Furthermore, our findings reveal that reusing a single threshold through all~$k$ attempts, either the~\EER threshold or the single attempt optimal threshold, leads to a significant reduction in the mechanism's success probability. This reduction indicates that the attacker is more likely to succeed in spoofing the biometric credential as the number of attempts increases.

% Previously, we noted that for each number of attempts~$k$, the optimal threshold is constant across all~$k$ attempts and increases with~$k$. Thus, we investigate how the optimal threshold selections are affected when we constrain the mean number of attempts, seeking to maximize the mechanism's success probability. We find that the optimal threshold increases with the number of attempts~$k$, and that the mechanism's success probability increases as the mean number of attempts increases. Nevertheless, system designers seek to minimize the mean number of a user's authentication attempts. 

% We explore the tradeoff between the mean number of attempts and the mechanism's success probability. First, we identify the optimal threshold selections that minimize the mean number of user authentication attempts, while requiring the mechanism's success probability to exceed~$P_{\mathrm{success}}^{\min}$. This analysis shows that the greater the required mechanism's success probability~$P_{\mathrm{success}}^{\min}$, the greater the mean number of user authentication attempts~$\mathbb{E}[N_U]$ required to achieve it. In addition, our findings reveal that a higher success probability constraint leads to a greater cumulative user success probability and a lower cumulative attacker success probability. This behavior indicates that the selected thresholds tend to be stricter in the initial attempts, with smaller differences between successive thresholds as the attempt index increases. In this way, the mechanism limits the probability of a successful spoofing attack, while relying on the user's cumulative success probability to increase with the number of attempts. These findigns are summarized in Figure~\ref{fig:success_constraints_comparison}.

% We also analyze the reverse direction, in which a system designer seeks to maximize the achievable security under the usability constraint. Specifically, we constrain the mean number of user authentication attempts~$\mathbb{E}[N_U]$, and seek to maximize the mechanism's success probability. We find that allowing a greater mean number of user authentication attempts~$\mathbb{E}[N_U]$, leads to a lower cumulative success probability for the attacker. 
% A higher limit on the mean number of attempts leads to stricter threshold selections in the initial authentication attempts. In this way, the mechanism prioritizes preventing the attacker from spoofing the biometric credential, over accepting the user in the initial attempts. Thus, the user's cumulative authentication probability is initially lower than under a lower limit on the mean number of attempts, but increases significantly over subsequent attempts as the user has more opportunities to obtain a score above the threshold.
% These findings are shown in Figure~\ref{fig:mean_limits_comparison}. 

% \begin{figure}[t]
% \centering

% % Upper subfigure: probabilities versus maximum attempt count
% \begin{subfigure}{\columnwidth}
% \centering
% \begin{tikzpicture}
% \begin{axis}[
%     width=0.95\linewidth,
%     height=0.65\linewidth,
%     xlabel={Number of attempts, $k$},
%     ylabel={Probability},
%     xmin=1,
%     xmax=10,
%     ymin=0,
%     ymax=1,
%     xtick=data,
%     grid=major,
%     line width=1pt,
%     mark size=1.6pt,
%     legend style={
%         at={(0.5,-0.28)},
%         anchor=north,
%         legend columns=3,
%         draw=black,
%         fill=white,
%         font=\scriptsize
%     }
% ]

% \addplot[green!60!black, solid, mark=*]
% table[x=k, y=P_success_max]
% {figs/fig_spoofed_users/livdet_fixed_prefix_far_frr_success_by_k.txt};
% \addlegendentry{$P_{\mathrm{success}}^{\max}$}

% \addplot[blue, dashed, mark=square*]
% table[x=k, y=P_success_EER]
% {figs/fig_spoofed_users/livdet_fixed_prefix_far_frr_success_by_k.txt};
% \addlegendentry{$P_{\mathrm{success}}(t_{\mathrm{EER}})$}

% \addplot[red, densely dotted, mark=triangle*]
% table[x=k, y=P_success_single_attempt_optimum]
% {figs/fig_spoofed_users/livdet_fixed_prefix_far_frr_success_by_k.txt};
% \addlegendentry{$P_{\mathrm{success}}(t_{\mathrm{opt}}^{(1)})$}

% \addplot[orange, solid, mark=diamond*]
% table[x=k, y=product_FRR]
% {figs/fig_spoofed_users/livdet_fixed_prefix_far_frr_success_by_k.txt};
% \addlegendentry{$\prod_{j=1}^{k}\mathrm{FRR}_j$}

% \addplot[purple, solid, mark=pentagon*]
% table[x=k, y=product_1_minus_FAR]
% {figs/fig_spoofed_users/livdet_fixed_prefix_far_frr_success_by_k.txt};
% \addlegendentry{$\prod_{j=1}^{k}(1-\mathrm{FAR}_j)$}

% \end{axis}
% \end{tikzpicture}
% \caption{Cumulative success and rejection probabilities.}
% \label{fig:Success_probability_by_k_probabilities}
% \end{subfigure}

% \medskip

% % Lower subfigure: final applied threshold for each k-attempt policy
% \begin{subfigure}{\columnwidth}
% \centering
% \begin{tikzpicture}
% \begin{axis}[
%     width=0.95\linewidth,
%     height=0.50\linewidth,
%     xlabel={Number of attempts, $k$},
%     ylabel={Selected threshold, $t_k^{(k)}$},
%     xmin=1,
%     xmax=10,
%     xtick=data,
%     grid=major,
%     line width=1pt,
%     mark size=1.6pt
% ]

% \addplot[green!60!black, solid, mark=*]
% table[x=k, y=Threshold]
% {figs/fig_spoofed_users/livdet_fixed_prefix_far_frr_success_by_k.txt};

% \end{axis}
% \end{tikzpicture}
% \caption{Threshold applied at the final attempt of each policy.}
% \label{fig:Success_probability_by_k_thresholds}
% \end{subfigure}

% \caption{Success probabilities and selected thresholds as a function
% of the maximum number of authentication attempts~$k$.
% (a)~The optimized cumulative success probability, compared with
% reusing the \EER threshold and the single-attempt optimal threshold,
% together with the products of user-rejection and attacker-rejection
% probabilities.
% (b)~The threshold applied at attempt~$k$ of the corresponding
% $k$-attempt fixed-prefix policy.
% Each value of~$k$ represents a separately optimized policy.}
% \label{fig:Success_probability_by_k}
% \end{figure}

% \pgfplotsset{compat=1.18}

% % Document
% \begin{figure}[t]
% \centering
% \begin{tikzpicture}
% \begin{axis}[
%     width=0.95\columnwidth,
%     height=0.65\columnwidth,
%     xlabel={Number of attempts, $k$},
%     ylabel={Overall success probability},
%     xmin=1, xmax=200,
%     ymin=0, ymax=1.02,
%     xtick={1,25,50,75,100,125,150,175,200},
%     ytick={0,0.2,0.4,0.6,0.8,1},
%     grid=major,
%     legend style={
%         at={(0.5,-0.31)},
%         anchor=north,
%         legend columns=1,
%         draw=black,
%         fill=white,
%         font=\scriptsize
%     },
%     legend cell align=left
% ]

% \addplot[
%     blue,
%     thick,
%     mark=o,
%     mark size=1.8pt,
%     mark repeat=10,
%     mark options={solid}
% ]
% table[
%     x=k,
%     y=P_success_optimized
% ] {figs/fig_spoofed_users/livdet_unconstrained_success_convergence.txt};
% \addlegendentry{Optimized overall success}

% \end{axis}
% \end{tikzpicture}
% \caption{Optimized mechanism success probability as a function of
% the number of authentication attempts~$k$.
% The threshold policy is optimized separately for each~$k$,
% without constraints on the conditional mean number of attempts
% or the success probability. The dashed line marks probability one.}
% \label{fig:unconstrained_success_convergence}
% \end{figure}

% \pgfplotsset{compat=1.18}

% \pgfplotsset{
%     constraint axis/.style={
%         width=0.9\linewidth,
%         height=0.55\linewidth,
%         xlabel={Attempt},
%         xmin=1, xmax=10,
%         xtick={1,...,10},
%         grid=major,
%         mark size=2pt,
%         legend style={
%             at={(0.5,-0.31)},
%             anchor=north,
%             draw=black,
%             fill=white,
%             font=\scriptsize
%         },
%         legend cell align=left
%     },
%     player axis/.style={
%         constraint axis,
%         ylabel={Cumulative success probability},
%         ymin=0, ymax=1,
%         legend columns=3,
%         every axis plot/.append style={forget plot}
%     }
% }

% % Two-row legend:
% % Row 1: three constraint values and their line styles.
% % Row 2: user and attacker markers.
% \newcommand{\TwoRowConstraintLegend}[3]{%
%     \addlegendimage{black, thick, solid, mark=none}
%     \addlegendentry{#1}
%     \addlegendimage{black, thick, densely dotted, mark=none}
%     \addlegendentry{#2}
%     \addlegendimage{black, thick, dashed, mark=none}
%     \addlegendentry{#3}
%     \addlegendimage{blue, only marks, mark=o, mark size=2pt}
%     \addlegendentry{User}
%     \addlegendimage{red, only marks, mark=square, mark size=2pt}
%     \addlegendentry{Attacker}
%     \addlegendimage{empty legend}
%     \addlegendentry{}
% }

% % =========================================================
% % Figure 1: minimize mean attempts under success constraints
% % =========================================================

% \begin{figure}[t]
% \centering
% \def\PolicyData{figs/fig_spoofed_users/livdet_success_constraints_comparison.txt}

% % Upper: threshold versus attempt
% \begin{subfigure}{\linewidth}
% \centering
% \begin{tikzpicture}
% \begin{axis}[
%     constraint axis,
%     ylabel={Threshold},
%     ymin=0, ymax=100
% ]

% \addplot[black, thick, solid, mark=*, mark size=1pt, mark options={solid}]
% table[x=Attempt, y=Threshold_0_6] {\PolicyData};

% \addplot[black, thick, densely dotted, mark=*, mark size=1pt, mark options={solid}]
% table[x=Attempt, y=Threshold_0_7] {\PolicyData};

% \addplot[black, thick, dashed, mark=*, mark size=1pt, mark options={solid}]
% table[x=Attempt, y=Threshold_0_75] {\PolicyData};

% \end{axis}
% \end{tikzpicture}
% \caption{Selected thresholds.}
% \label{fig:success_constraints_thresholds}
% \end{subfigure}

% \par\medskip

% % Lower: user and attacker success versus attempt
% \begin{subfigure}{\linewidth}
% \centering
% \begin{tikzpicture}
% \begin{axis}[player axis]

% % Lowest constraint: solid
% \addplot[blue, thick, solid, mark=o, mark options={solid}]
% table[x=Attempt, y=P_user_success_0_6] {\PolicyData};

% \addplot[red, thick, solid, mark=square, mark options={solid}]
% table[x=Attempt, y=P_attacker_success_0_6] {\PolicyData};

% % Middle constraint: dotted
% \addplot[blue, thick, densely dotted, mark=o, mark options={solid}]
% table[x=Attempt, y=P_user_success_0_7] {\PolicyData};

% \addplot[red, thick, densely dotted, mark=square, mark options={solid}]
% table[x=Attempt, y=P_attacker_success_0_7] {\PolicyData};

% % Highest constraint: dashed
% \addplot[blue, thick, dashed, mark=o, mark options={solid}]
% table[x=Attempt, y=P_user_success_0_75] {\PolicyData};

% \addplot[red, thick, dashed, mark=square, mark options={solid}]
% table[x=Attempt, y=P_attacker_success_0_75] {\PolicyData};

% \TwoRowConstraintLegend
%     {$P'=0.6$}
%     {$P'=0.7$}
%     {$P'=0.75$}

% \end{axis}
% \end{tikzpicture}
% \caption{Cumulative user and attacker acceptance probabilities.}
% \label{fig:success_constraints_probabilities}
% \end{subfigure}

% \caption{Policies minimizing the conditional mean number of user
% authentication attempts subject to
% $P_{\mathrm{success}}\geq P'$, with $k=10$.
% The upper panel shows selected thresholds; the lower panel shows
% the probability of at least one user or attacker acceptance through
% each attempt. Line styles identify constraint values, while blue
% circles represent the user and red squares represent the attacker.
% The constraint applies to the complete policy.}
% \label{fig:success_constraints_comparison}
% \end{figure}

% % =========================================================
% % Figure 2: maximize success under mean-attempt constraints
% % =========================================================

% \begin{figure}[t]
% \centering
% \def\PolicyData{figs/fig_spoofed_users/livdet_mean_limits_comparison.txt}

% % Upper: threshold versus attempt
% \begin{subfigure}{\linewidth}
% \centering
% \begin{tikzpicture}
% \begin{axis}[
%     constraint axis,
%     ylabel={Threshold},
%     ymin=0, ymax=100
% ]

% \addplot[black, thick, solid, mark=*, mark size=1pt, mark options={solid}]
% table[x=Attempt, y=Threshold_1_5] {\PolicyData};

% \addplot[black, thick, densely dotted, mark=*, mark size=1pt, mark options={solid}]
% table[x=Attempt, y=Threshold_1_9] {\PolicyData};

% \addplot[black, thick, dashed, mark=*, mark size=1pt, mark options={solid}]
% table[x=Attempt, y=Threshold_2_4] {\PolicyData};

% \end{axis}
% \end{tikzpicture}
% \caption{Selected thresholds.}
% \label{fig:mean_limits_thresholds}
% \end{subfigure}

% \par\medskip

% % Lower: user and attacker success versus attempt
% \begin{subfigure}{\linewidth}
% \centering
% \begin{tikzpicture}
% \begin{axis}[player axis]

% % Lowest constraint: solid
% \addplot[blue, thick, solid, mark=o, mark options={solid}]
% table[x=Attempt, y=P_user_success_1_5] {\PolicyData};

% \addplot[red, thick, solid, mark=square, mark options={solid}]
% table[x=Attempt, y=P_attacker_success_1_5] {\PolicyData};

% % Middle constraint: dotted
% \addplot[blue, thick, densely dotted, mark=o, mark options={solid}]
% table[x=Attempt, y=P_user_success_1_9] {\PolicyData};

% \addplot[red, thick, densely dotted, mark=square, mark options={solid}]
% table[x=Attempt, y=P_attacker_success_1_9] {\PolicyData};

% % Highest constraint: dashed
% \addplot[blue, thick, dashed, mark=o, mark options={solid}]
% table[x=Attempt, y=P_user_success_2_4] {\PolicyData};

% \addplot[red, thick, dashed, mark=square, mark options={solid}]
% table[x=Attempt, y=P_attacker_success_2_4] {\PolicyData};

% \TwoRowConstraintLegend
%     {$\mu'=1.5$}
%     {$\mu'=1.9$}
%     {$\mu'=2.4$}

% \end{axis}
% \end{tikzpicture}
% \caption{Cumulative user and attacker acceptance probabilities.}
% \label{fig:mean_limits_probabilities}
% \end{subfigure}

% \caption{Policies maximizing the mechanism's success probability
% subject to $\mathbb{E}[N_U\mid N_U\leq k]\leq\mu'$,
% with $k=10$.
% The upper panel shows selected thresholds; the lower panel shows
% the probability of at least one user or attacker acceptance through
% each attempt. Line styles identify constraint values, while blue
% circles represent the user and red squares represent the attacker.
% The constraint applies to the complete policy.}
% \label{fig:mean_limits_comparison}
% \end{figure}

\section{Biometric Sensors Survey}
\label{sec:sensors_review}
\noindent
We review manufacturers' specifications of commercially available biometric sensors, spanning fingerprints, face, iris and palm.
We identify 32 device families from 21 manufacturers and classify each family according to three criteria that reflect our model's requirements.

The first criterion is \emph{in-sensor decision}: whether the device performs matching and emits an accept/reject verdict itself, or merely captures a sample for a host matcher.
Of the 32 families, 23 make an in-sensor decision while the remaining 9 are capture only.
A capture-only sensor exposes the biometric sample or the matching score, which can facilitate hill-climbing attacks~\cite{Bringer2018PrivateByDesign} and creates privacy risks, as third parties may reuse the samples to identify users across services~\cite{FIDOPrivacyPrinciples}.

The second criterion is \emph{threshold tuning}: whether the device allows the integrator to set a numeric match threshold.
Only 2 of the 23 in-sensor decision families do, comprising 6 individual models.
The remaining families offer weaker methods that do not permit setting a specific threshold: a target error rate within a predetermined range, a discrete security level such as~1--5 or High/Medium/Low, or no threshold parameter at all.
In each case the integrator is left with a threshold the sensor derives internally and never exposes.

The third criterion is \emph{error-rate exposure}: whether the manufacturer publishes the~\FAR and~\FRR functions for the device.
Setting a threshold is necessary but not sufficient for optimization, since without these functions the integrator cannot determine the error rates a threshold induces.
None of the 23 families satisfies this criterion.
The detailed list is in Appendix~\ref{app:Bio_table}.

The practical consequence is that today's authentication mechanisms cannot optimize the match threshold they inherit.
In 21 of the 23 families we surveyed the threshold cannot be set at all, and in the remaining two it can be set but not informed, since the error-rate functions are not published.
To support threshold optimization, manufacturers should expose both the match threshold and the~\FAR and~\FRR functions that characterize it.
This would allow integrators to select thresholds according to their requirements, and as part of the other credentials used in the mechanism.

\section{Conclusion}\label{sec:conclusion}
\noindent
In a single-attempt mechanism using a biometric credential alone, optimizing the threshold yields only small improvements, which may explain why it has received limited attention. But when the credential is combined with other credentials or used across multiple attempts, the optimal threshold depends on the mechanism and differs from the heuristic one. We characterize a condition under which adding a biometric credential always improves the mechanism’s success, and prove it. We also find that the optimal threshold depends on the number of allowed attempts, and present two constrained optimization problems that let system designers balance security and usability. Simulating spoofing attacks on public biometric datasets, we find that threshold optimization reduces the failure probability by~20-72\% compared to the heuristic threshold, and by up to~21\% compared to using the atomic credential alone. Reusing the heuristic threshold across attempts reduces the success probability by up to 72 percentage points.

However, our survey of 32 commercial device families from 21 manufacturers finds that only two allow setting a numeric match threshold, and none publishes the~\FAR and~\FRR functions required to choose one. We therefore propose that manufacturers should expose both, so integrators can select thresholds for their mechanisms. This paves the way for developing new authentication protocols based on threshold selection and may enhance the security of digital and physical assets for millions of users around the world.

{
\raggedright
\emergencystretch=3em
\Urlmuskip=0mu plus 2mu\relax
\def\UrlBreaks{%
  \do\a\do\b\do\c\do\d\do\e\do\f\do\g\do\h\do\i\do\j\do\k\do\l\do\m
  \do\n\do\o\do\p\do\q\do\r\do\s\do\t\do\u\do\v\do\w\do\x\do\y\do\z
  \do\A\do\B\do\C\do\D\do\E\do\F\do\G\do\H\do\I\do\J\do\K\do\L\do\M
  \do\N\do\O\do\P\do\Q\do\R\do\S\do\T\do\U\do\V\do\W\do\X\do\Y\do\Z
  \do\0\do\1\do\2\do\3\do\4\do\5\do\6\do\7\do\8\do\9
  \do\/\do\-\do\.\do\_\do\?\do\&\do\=\do\#}
\bibliographystyle{IEEEtran}
\bibliography{main}
}

\appendices

\section{A Sufficient Condition for Integrating Biometric Credentials}
\label{app:sufficient-condition}
\noindent
Since biometric credentials often exhibit a more common behavior~\cite{Slob2005,anil2004}, we introduce an \emph{asymptotic dominancy} condition which captures this property. We prove that adding an \emph{asymptotically dominated} biometric credential can produce a new mechanism~$\bar{M}$ with higher success probability than that of a single atomic credential mechanism~$M$, and show it is a sufficient condition.
\begin{definition}
\label{def:User and Attacker distribution functions}
A biometric credential with a user distribution function~$f_U(s)$ and an attacker distribution function~$f_A(s)$ are \emph{asymptotically dominated} if
\begin{equation*}
\begin{aligned}
&\forall c>0\;\exists t\; \text{s.t.}:\forall s \ge t: f_A(s)\le c\,f_U(s),\textit{ or} \\
&\forall c>0\;\exists t\; \text{s.t.}:\forall s \le t: f_U(s)\le c\,f_A(s).
\end{aligned}
\end{equation*}
\end{definition}

\begin{proposition}\label{thm:biometric-improves-security}
Let~$c_1$ be an atomic credential with~$P^1_\text{safe} > 0$ and at least one of~$P^1_\text{leak},P^1_\text{loss} > 0$, $c_2$ be a biometric credential whose user distribution~$f_U(t)$ and attacker distribution~$f_A(t)$ are asymptotically dominated and~$M$ be a mechanism using~$c_1$. Then, there exists a mechanism~$\bar{M}$ using~$c_1$,~$c_2$, such that
    \begin{equation*}
        P_{\mathrm{success}}^{\bar{M}} > P_{\mathrm{success}}^{M}
    \end{equation*}
\end{proposition}

\begin{proof}
The success probability of a mechanism comprising an atomic credential only is~$P^{1}_\textit{safe}$~(Equation~\ref{eqn:safesuccess}). A biometric credential is characterized by four state probabilities that are functions of~\FAR and~\FRR
\contprobabilities
We previously observed that a sufficient condition for producing a new mechanism with greater success is described by Equation~\ref{eqn:FRRineq} and Equation~\ref{eqn:FARineq}. Thus, we need to show that there exists an operating threshold $t$ such that either
\begin{equation*}
    \begin{aligned}
    &\frac{\FRR(t)}{1-\FAR(t)} < \frac{P^{1}_{\text{leak}}}{P^{1}_{\text{safe}}+P^{1}_{\text{leak}}} \textit{, or } \\
    &\frac{\FAR(t)}{1-\FRR(t)} < \frac{P^{1}_{\text{loss}}}{P^{1}_{\text{safe}}+P^{1}_{\text{loss}}}
    \end{aligned}
\end{equation*}
We expand left hand side of Equation~\ref{eqn:FRRineq}. 
\begin{equation*}
\frac{\FRR(t)}{1-\FAR(t)} = \frac{\int_{-\infty}^{t} f_U(s) ds}{\int_{-\infty}^{t} f_A(s) ds}
\end{equation*}
By Definition~\ref{def:User and Attacker distribution functions}, since the biometric credential is asymptotically dominated, it follows that there exists an operating threshold~$t$ such that for every~$s \le t$, $f_U(s) \le c\cdot\,f_A(s)$. Therefore:
\begin{equation*}
\frac{\int_{-\infty}^{t} f_U(s) ds}{\int_{-\infty}^{t} f_A(s) ds} \le \frac{\int_{-\infty}^{t} c \cdot f_A(s) ds}{\int_{-\infty}^{t} f_A(s) ds} = c
\end{equation*}
This holds~$\forall c > 0,$; thus, we can choose~$c$ to equal the right hand side of~Equation~\ref{eqn:FRRineq} and satisfy the sufficient condition.
Similarly, we expand left hand side of Equation~\ref{eqn:FARineq}. 
\begin{equation*}
\frac{\FAR(t)}{1-\FRR(t)} = \frac{\int_{t}^{\infty} f_A(s) ds}{\int_{t}^{\infty} f_U(s) ds}
\end{equation*}
By Definition~\ref{def:User and Attacker distribution functions}, since the biometric credential is asymptotically dominated, it follows that there exists an operating threshold~$t$ such that for every~$s \ge t,$ $f_A(s) \le c\cdot\,f_U(s)$. Therefore,
\begin{equation*}
\frac{\int_{t}^{\infty} f_A(s) ds}{\int_{t}^{\infty} f_U(s) ds} \le \frac{\int_{t}^{\infty} c \cdot f_U(s) ds}{\int_{t}^{\infty} f_U(s) ds} = c, \forall c > 0
\end{equation*}
We have showed that~$\forall c>0$ either~$\frac{\FAR(t)}{1-\FRR(t)} < c$ or $\frac{\FRR(t)}{1-\FAR(t)} < c$. Therefore, integrating an asymptotically dominated biometric credential with an atomic credential produces~$\bar{M}$ with greater success probability.
\end{proof}

This proof is non trivial, as the integration of an atomic credential mechanism with another atomic credential, does not necessarily produce a new mechanism with higher success probability. Moreover, we have shown that integrating a biometric credential into an atomic credential mechanism without adjusting the match threshold~$t$ may reduce the mechanism's success probability~$(\S\ref{sec:atm_biocred})$. Thus, Proposition~\ref{thm:suf_nec_cond_for_improve} and Proposition~\ref{thm:biometric-improves-security} emphasize the importance of adjusting the threshold when adding a biometric credential into an authentication mechanism.
The \emph{asymptotically dominated} condition reflects practical scenarios, as for very low score values~$s$ the attacker distribution dominates the user distribution, and vice versa for very high values of~$s$. We validate this claim with real public biometric data in Section~\ref{sec:empirical_evaluation}.

\section{Biometric Sensors Survey Table}\label{app:Bio_table}
\noindent
Table~\ref{tab:device_survey} shows all the commercially available devices we surveyed. Rows are grouped by threshold exposure, from devices permitting an arbitrary match threshold down to those exposing none, followed by the capture-only devices outside our scope. The~\emph{On-sensor} column indicates whether the device performs matching and emits an accept/reject verdict itself. The~\emph{Threshold exposure} column gives the mechanism, if any, by which an integrator may influence the decision boundary. Families from one manufacturer sharing a modality and threshold mechanism are listed on a single row, so the family counts in the group headings exceed the number of rows. Each entry cites the manufacturer documentation from which its classification was drawn.
% ---------------------------------------------------------------------------
% Comparison of surveyed biometric device families.
% Requires: booktabs only. Spans both columns (table*).
% Families sharing a manufacturer, modality and threshold mechanism share a row;
% group headings give the family count, which exceeds the row count.
% ---------------------------------------------------------------------------

\begin{table*}[t]
\centering
\caption{Surveyed biometric device families, grouped by threshold exposure.
On-sensor column indicates the device performs matching and emits an accept/reject
verdict itself. Families from one manufacturer sharing a modality and threshold
mechanism are listed on a single row.}
\label{tab:device_survey}
\footnotesize
\setlength{\tabcolsep}{4pt}
\begin{tabular}{@{}llllll@{}}
\toprule
\textbf{Device family} & \textbf{Manufacturer} & \textbf{Modality} & \textbf{On-sensor} & \textbf{Threshold exposure} & \textbf{Ref.} \\
\midrule
\multicolumn{6}{@{}l}{\textit{Arbitrary match threshold (2 families, 6 models)}} \\
\midrule
{F18}, {SpeedFace-V5L}, {MultiBio700}, {SF100}, {F21} & ZKTeco & FP / face / palm & Yes & Numeric score & \cite{zkteco_threshold_manual,zkteco_speedface_v5l,zkteco_terminal_manual} \\
{DS-K1T342} series & Hikvision & Face / FP & Yes$^{\dagger}$ & Numeric score & \cite{hikvision_dsk1t342,hikvision_isapi} \\
\midrule
\multicolumn{6}{@{}l}{\textit{Target error rate (2 families)}} \\
\midrule
{EF-45}, {EF-45N}, {EF-45NC}, {EF-45}~{AC} & CMITech & Iris / face & Yes & {FMR} $10^{-5}$--$10^{-14}$ & \cite{cmitech_ef45} \\
{MorphoAccess} {SIGMA} Lite/Lite+/Wide/Extreme & IDEMIA & FP / face & Yes & {FAR} 1\%--$10^{-7}$\% & \cite{idemia_morphoaccess_sigma,idemia_sigma_lite} \\
\midrule
\multicolumn{6}{@{}l}{\textit{Discrete security level (12 families)}} \\
\midrule
{R307} / {AS608}, {R503}, {R502-A} & Hangzhou Grow & FP & Yes & Level 1--5 & \cite{grow_r307,grow_r503,grow_r502a} \\
{ID809} ({SEN0348} / {SEN0359}) & DFRobot & FP & Yes & Level 1--5 & \cite{dfrobot_id809} \\
{NAC-3000}, Fingkey Access & NITGEN & FP & Yes & Level 1--9, per mode & \cite{nitgen_nac3000} \\
{FIM5360} / {FIM6060} / {FIM5260} / {FIM4060} & NITGEN & FP & Yes & Level, per user & \cite{nitgen_fim5360} \\
{AC-7000} & Union Community & Face / FP & Yes & Level 1--9, per mode & \cite{virdi_ac7000} \\
{ID} {F450} / {F455} & RealSense & Face & Yes & High/Med/Low & \cite{realsense_id} \\
{BioStation}~3, {FaceStation}~{F2} & Suprema & Face / FP & Yes & Level, not enumerated & \cite{suprema_biostation3,suprema_facestationf2} \\
{SFM} Slim, {SFM5500} & Suprema & FP & Yes & {FAR} values, per user & \cite{suprema_uf_protocol,suprema_sfm5500} \\
\midrule
\multicolumn{6}{@{}l}{\textit{No threshold parameter (7 families)}} \\
\midrule
{GT-521F32} / {GT-521F52} & ADH-Tech & FP & Yes & None & \cite{adhtech_gt521f} \\
{BM-Lite} ({SafeTouch}) & Fingerprint Cards & FP & Yes & None & \cite{fpc_bmlite} \\
{IriShield} {MK}~2120{U} / {UL} & IriTech & Iris & Yes & Not documented & \cite{iritech_irishield} \\
{FVA-U4ST} & mofiria & Finger vein & Yes & Not documented & \cite{mofiria_fvau4st} \\
{iCAM7000} series & Iris ID & Iris & Mode-dep.$^{\ddagger}$ & Not documented & \cite{irisid_icam7000} \\
{PalmSecure} {ID} Match & Fujitsu & Palm vein & Yes & Not documented & \cite{fujitsu_idmatch} \\
{IXM} {TITAN} / {TOUCH}~2 / {TFACE} & Invixium & Face / FP / vein & Yes & Unverified & \cite{invixium_ixm} \\
\midrule
\multicolumn{6}{@{}l}{\textit{Capture only --- outside scope (9 families)}} \\
\midrule
{DigitalPersona} 4500 & HID Global & FP & No & Host: numeric & \cite{hid_digitalpersona4500,hid_uareu_sdk} \\
Hamster Pro 20 ({HU20}) & SecuGen & FP & No & Host: score 0--199 & \cite{secugen_hamsterpro20,secugen_fdx_sdk} \\
{BioMini} Slim~2 & Suprema & FP & No & Host {SDK} & \cite{suprema_biominislim2} \\
{ZK9500} / {SLK20R} & ZKTeco & FP & No & Host {SDK} & \cite{zkteco_slk20r} \\
{MFS100} & Mantra Softech & FP & No & Host {SDK} / {RD} service & \cite{mantra_mfs100} \\
{PalmSecure}-{F} {Pro}+ & Fujitsu & Palm vein & No & Host library & \cite{fujitsu_fproplus} \\
{H-1} ({PC-KCA100}/{110}) & Hitachi & Finger vein & No & Host engine & \cite{hitachi_h1} \\
{BMT-20} & CMITech & Iris & No & Host {SDK} & \cite{cmitech_bmt20} \\
{SEN0626} & DFRobot & Face & No matcher & n/a$^{\S}$ & \cite{dfrobot_sen0626} \\
\bottomrule
\end{tabular}

\vspace{2pt}
\footnotesize
$^{\dagger}$Self-deciding only in Terminal Authentication Mode; in Reader
Authentication Mode the decision is delegated to an upstream controller.
$^{\ddagger}$On-device matching is one of three selectable operating modes.
$^{\S}$Face detection only: no template store and no identity decision.
\end{table*}

\section{Closed Form of the Quadratic Optimal Threshold}\label{app:quad_opt_thr}
\noindent
We derive the optimal threshold for a quadratic distribution, which is a second-order approximation of a biometric credential score distribution. As shown in~$(\S\ref{sec:quadratic_biocred})$, the optimal threshold is the single real root of the cubic in the third interval.
Let
\begin{align*}
p &= 13\alpha_1^2 - 2\alpha_1(13\alpha_2 + \beta_1 - \beta_2) + 25\alpha_2^2 \\
  &\quad - 2\alpha_2(11\beta_1 + \beta_2) + 25\beta_1^2 - 26\beta_1\beta_2 + 13\beta_2^2, \\[4pt]
r &= 19\alpha_1^2 + 2\alpha_1(35\alpha_2 - 55\beta_1 + \beta_2) - 125\alpha_2^2 \\
  &\quad + 10\alpha_2(29\beta_1 - 11\beta_2) - (25\beta_1 - 19\beta_2)(5\beta_1 + \beta_2), \\[4pt]
w &= 19\alpha_1^3 + 3\alpha_1^2(17\alpha_2 - 43\beta_1 + 7\beta_2) \\
  &\quad + 3\alpha_1\bigl(-5(13\alpha_2^2 - 22\alpha_2\beta_1 + \beta_1^2) \\
  &\quad\quad - 14\beta_2(\alpha_2 + \beta_1) + 7\beta_2^2\bigr) \\
  &\quad + 125\alpha_2^3 - 165\alpha_2^2\beta_1 - 15\alpha_2^2\beta_2 - 165\alpha_2\beta_1^2 \\
  &\quad + 330\alpha_2\beta_1\beta_2 - 129\alpha_2\beta_2^2 + 125\beta_1^3 \\
  &\quad - 195\beta_1^2\beta_2 + 51\beta_1\beta_2^2 + 19\beta_2^3, \\[4pt]
D &= \sqrt{-p^3 + (\alpha_1 - \alpha_2 - \beta_1 + \beta_2)^2 r^2}, \\[4pt]
C &= \sqrt[3]{w + D}.
\end{align*}
The optimal threshold is then
\begin{equation*}
t_{\text{opt}} = \frac{1}{12}\left(C + \frac{p}{C} + \alpha_1 + 5\alpha_2 + 5\beta_1 + \beta_2\right).
\end{equation*}

\end{document}
\endinput
